\exercisegroup

\begin{enumerate}
\def\labelenumi{\arabic{enumi})}
\tightlist
\item
  \begin{enumerate}
  \def\labelenumii{\alph{enumii})}
  \tightlist
  \item
    Let \(\mathbf{E}\) be a locally convex Hausdorff space, and \(\mathbf{E}'\) its dual. Show that the following properties are equivalent:
  \end{enumerate}
\end{enumerate}

α) E is infra-barrelled (III, p.~44, exerc. 7).

β) Every strongly bounded subset in E' is equicontinuous.

\(\gamma)\) Every strongly bounded subset in \(E'\) is relatively weakly compact, and the initial topology of \(E\) is \(\tau(E, E')\) .

δ) The topology induced on E by the strong topology of the bidual \(E''\) is identical with the initial topology on E.

Then a fundamental system of neighbourhoods of 0 for the strong topology of \(\mathbf{E}''\) consists of the closures, for the topology \(\sigma (\mathrm{E}'', \mathrm{E}')\), of a fundamental system of neighbourhoods of 0 for the initial topology of \(\mathbf{E}\) .

\begin{enumerate}
\def\labelenumi{\alph{enumi})}
\setcounter{enumi}{1}
\tightlist
\item
  Prove that if E is infra-barrelled and if its dual \(E'\) is identical with its algebraic dual \(E^{*}\) , then the initial topology of E is the finest locally convex topology.
\end{enumerate}

¶ 2) a) Show that every product of infra-barrelled spaces is infra-barrelled. (Reduce to the case of Hausdorff infra-barrelled spaces; then use exerc. 1 and prop. 15 of IV, p.~14.)

\begin{enumerate}
\def\labelenumi{\alph{enumi})}
\setcounter{enumi}{1}
\tightlist
\item
  Give an example of a Hausdorff and infra-barrelled locally convex space which is neither bornological non barrelled. (Proceed as in III, p.~45, exerc. 16, replacing the barrelled spaces by infra-barrelled spaces and use a).)
\end{enumerate}

\begin{enumerate}
\def\labelenumi{\arabic{enumi})}
\setcounter{enumi}{2}
\item
  Let E be a complex Hausdorff locally convex space, \(E_{0}\) the underlying real locally convex space to E, and let \(E^{\prime}\) and \(E_{0}^{\prime}\) be the duals of E and \(E_{0}\) respectively. Show that the canonical R-linear mapping \(f \mapsto Rf\) from \(E^{\prime}\) onto \(E_{0}^{\prime}\) is a homeomorphism for the S-topologies on \(E^{\prime}\) and \(E_{0}^{\prime}\) , where S is an arbitrary set of bounded subsets of E. From this deduce the definition of the canonical R-linear mapping from the bidual \(E^{\prime\prime}\) onto the bidual \(E_{0}^{\prime\prime}\) , which is a homeomorphism for the weak topologies \(\sigma(E^{\prime\prime}, E^{\prime})\) and \(\sigma(E_{0}^{\prime\prime}, E_{0}^{\prime})\) , as also for the strong topologies \(\beta(E^{\prime\prime}, E^{\prime})\) and \(\beta(E_{0}^{\prime\prime}, E_{0}^{\prime})\) ; by this mapping, E (considered embedded in \(E^{\prime\prime}\) ) transforms into \(E_{0}\) (considered embedded in \(E_{0}^{\prime\prime}\) ).
\item
  Let E be a Hausdorff and infra-barrelled locally convex space.
\end{enumerate}

\begin{enumerate}
\def\labelenumi{\alph{enumi})}
\tightlist
\item
  Show that if the strong dual \(E_{b}^{\prime}\) of E is bornological, then the completion \(\hat{E}\) of E, identified with a vector subspace of \(E^{\prime*}\) (III, p.~21, th. 2), is contained in the bidual \(E^{\prime\prime}\) of E.
\item
  We say that E is distinguished if every subset of \(E^{\prime\prime}\) which is bounded for the topology \(\sigma(\mathrm{E}^{\prime\prime}, \mathrm{E})\) is contained in the closure (for this topology) of a bounded subset of E. Show that for E to be distinguished it is necessary and sufficient that its strong dual \(E_{b}^{\prime}\) is barrelled (cf.~IV, p.~15, prop. 3).
\end{enumerate}

\begin{enumerate}
\def\labelenumi{\arabic{enumi})}
\setcounter{enumi}{4}
\tightlist
\item
  Let \(\mathbf{E}\) be a Hausdorff locally convex space, and \(\mathbf{E}'\) its dual.
\end{enumerate}

\begin{enumerate}
\def\labelenumi{\alph{enumi})}
\item
  In order that the strong topology on \(\mathbf{E}'\) be identical with \(\tau (\mathrm{E}',\mathrm{E})\) , it is necessary and sufficient that \(\mathbf{E}\) is semi-reflexive.
\item
  Suppose that E is infra-barrelled. In order that the strong topology on \(E'\) be identical with the topology of compact convergence, it is necessary and sufficient that E be a Montel space.
\end{enumerate}

¶ 6) Let F and G be two vector spaces in separating duality.

\begin{enumerate}
\def\labelenumi{\alph{enumi})}
\tightlist
\item
  Show that the following properties are equivalent :
\end{enumerate}

\(\alpha)\) the space \(\mathbf{F}\) with a topology compatible with the duality between \(\mathbf{F}\) and \(\mathbf{G}\) , is semi-reflexive;

\(\beta)\) the space \(\mathbf{G}\) , with \(\tau (\mathbf{G},\mathbf{F})\) , is barrelled.

\begin{enumerate}
\def\labelenumi{\alph{enumi})}
\setcounter{enumi}{1}
\tightlist
\item
  Show that the following properties are equivalent :
\end{enumerate}

α) the space F, with τ(F, G), is reflexive;

\(\beta)\) the space \(\mathbf{G}\) , with \(\tau (\mathbf{G},\mathbf{F})\) , is reflexive;

γ) F and G are barrelled for τ(F, G) and τ(G, F) respectively.

\begin{enumerate}
\def\labelenumi{\alph{enumi})}
\setcounter{enumi}{2}
\item
  Show that every locally convex Hausdorff space E, with the topology \(\tau(\mathrm{E}, \mathrm{E}')\) is isomorphic to the quotient of a semi-reflexive space F by a closed subspace. (By III, p.~44, exerc. 14, c), \(E'\) endowed with \(\tau(\mathrm{E}', \mathrm{E})\) is isomorphic to a closed vector subspace M of a Hausdorff barrelled space G; take \(F = G'\) endowed with \(\tau(\mathrm{G}', \mathrm{G})\) and use prop. 11 of IV, p.~10.)
\item
  Let A be an infinite set with \(\operatorname{Card}(\mathrm{A}) > \aleph_{1}\) (S, III, §6, exerc. 10). In the product space \(P = R^{A}\) , let \(E_{0}\) denote the everywhere dense subspace consisting of all \(x = (x_{\alpha})_{\alpha \in \mathrm{A}}\) such that \(x_{\alpha} = 0\) except for a countable set of indices; the space \(E_{0}\) is barrelled and so is the subspace E of P generated by \(E_{0}\) and the point \(1_{A}\) in P all whose coordinates are equal to 1 (III, p.~45, exerc. 16). Let B be a bounded subset in \(E_{0}\) , whose closure in P contains \(1_{A}\) , and let J be a subset of A with cardinality \(\aleph_{1}\) ; let \(pr_{J}\) denote the projection from P onto \(R^{J}\) ; show that there exists a subset \(B_{J}\) in B, with cardinality \(\leqslant \aleph_{1}\) , such that the closure of \(\operatorname{pr}_{J}(B_{J})\) in \(R^{J}\) contains \(1_{J}\) ; then the set \(J' \supset J\) of all indices \(\alpha \in A\) such that for at least one point of \(B_{J}\) the coordinate with index \(\alpha\) is \(\neq 0\) , has a cardinality equal to \(\aleph_{1}\) . We define \(J_{0} = J\) and by induction \(J_{n+1} = J'_{n}\) , and put \(H = \bigcup_{n} J_{n}\) , whose cardinality is \(\aleph_{1}\) , and let \(B_{H} = \bigcup_{n} B_{J_{n}}\) ; show that the closure of \(B_{H}\) (and hence also that of B) contains the point \((1_{H}, 0) \in R^{H} \times R^{A-H} = P\) , which does not belong to \(E_{0}\) . Conclude that for every bounded and compact set C in E, \(C \cap E_{0}\) is again closed in E, and that E is not semi-reflexive.
\item
  Deduce from \(d\) that the dual \(\mathrm{E}'\) of \(\mathrm{E}\) (which can be identified with the dual \(\mathrm{P}' = \mathbf{R}^{(\mathbf{A})}\) of \(\mathrm{P}\) ) is not complete for the topology \(\tau (\mathrm{E}',\mathrm{E})\) , in spite of being semi-reflexive (hence quasi-complete) for this topology (consider the linear form on \(\mathrm{E}\) which is equal to 0 on \(\mathrm{E}_0\) and equal to 1 at the point \(1_{\mathrm{A}}\) , and use III, p.~21, th. 2).
\end{enumerate}

\begin{enumerate}
\def\labelenumi{\arabic{enumi})}
\setcounter{enumi}{6}
\item
  Let \((\mathrm{E}_{i})_{i\in\mathbb{I}}\) be a family of Hausdorff locally convex spaces, P the product space of the \(E_{i}\) , and S their topological direct sum. Show that for P or S to be semi-reflexive (resp. reflexive), it is necessary and sufficient that each \(E_{i}\) be semi-reflexive (resp. reflexive).
\item
  Let E be a Hausdorff locally convex space which is the strict inductive limit of an increasing sequence \((E_{n})\) of closed vector subspaces (II, p.~32, prop. 9).
\end{enumerate}

\begin{enumerate}
\def\labelenumi{\alph{enumi})}
\item
  Show that, if the strong dual of each of the \(\mathbf{E}_n\) is complete, then the strong dual of \(\mathbf{E}\) is complete (III, p.~20, th. 1).
\item
  In order that E be semi-reflexive (resp. reflexive), it is necessary and sufficient that each of the \(E_{n}\) be semi-reflexive (resp. reflexive).
\end{enumerate}

\begin{enumerate}
\def\labelenumi{\arabic{enumi})}
\setcounter{enumi}{8}
\item
  Let \((\mathrm{E}_{\alpha})_{\alpha\in\mathrm{A}}\) be a family of Hausdorff locally convex spaces contained in the same vector space, which is directed for the relation \(\supset\) , such that, if \(E_{\beta} \subset E_{\alpha}\) , the topology of \(E_{\beta}\) is finer than the topology on \(E_{\beta}\) induced by that of \(E_{\alpha}\) . Let E be the intersection of the \(E_{\alpha}\) , endowed with a topology which is the supremum of the topologies on E induced by those on the \(E_{\alpha}\) . Show that if each \(E_{\alpha}\) is semi-reflexive, then E is semi-reflexive (consider an ultrafilter on a bounded subset of E).
\item
  Show that every product, and every topological direct sum of Montel spaces is a Montel space \(^{1}\) .
\item
  Let E be a Hausdorff locally convex space such that the strong dual \(E_{b}^{\prime}\) of E is semi-reflexive. a) Show that, on every strongly bounded subset of \(E^{\prime}\) , the topologies induced by \(\sigma(E^{\prime}, E^{\prime\prime})\) and \(\sigma(E^{\prime}, E)\) are identical.
\end{enumerate}

\begin{enumerate}
\def\labelenumi{\alph{enumi})}
\setcounter{enumi}{1}
\tightlist
\item
  Deduce from a) that E is infra-barrelled for the topology \(\tau(E, E')\) (cf.~IV, p.~52, exerc. 1) and that, if \(\hat{E}\) is its completion for this topology, and is identified with a subset of \(E'^{*}\) (III, p.~21, th. 2), then \(E'' \subset \hat{E}\) . In particular, if E is quasi-complete for \(\tau(E, E')\) , then E is reflexive for this topology.
\end{enumerate}

\(^{1}\) On the other hand, a closed subspace of a Montel space need not be infrabarrelled, and the quotient of a Montel space by a closed subspace need not be semi-reflexive (IV, p.~63, exerc. 8).

\begin{enumerate}
\def\labelenumi{\arabic{enumi})}
\setcounter{enumi}{11}
\tightlist
\item
  Let \(\mathbf{E}\) be a Banach space and \(\mathbf{E}'\) its dual.
\end{enumerate}

\begin{enumerate}
\def\labelenumi{\alph{enumi})}
\item
  Show that the distance \(x \mapsto d(x, \mathbf{A})\) from a point \(x \in \mathbf{E}\) to a closed convex set \(\mathbf{A}\) is a lower semi-continuous function on \(\mathbf{E}\) for the topology \(\sigma(\mathbf{E}, \mathbf{E}')\) .
\item
  Show that if \(\mathrm{E}\) is reflexive, then for every closed convex subset \(\mathrm{A}\) of \(\mathrm{E}\) , there exists a point \(x_0 \in \mathrm{A}\) such that \(\| x_0 \|\) is equal to the distance of 0 to \(\mathrm{A}\) (use \(a\) ). This point is unique if every boundary point of the unit ball of \(\mathrm{E}\) is extremal (II, p.~54, def. 1).
\item
  Suppose E is reflexive and let B be a closed convex and bounded subset in E; deduce from a) and b) that there exist two points \(x \in \mathbf{A}\) , \(y \in \mathbf{B}\) such that \(\| x - y \| = d(\mathbf{A}, \mathbf{B})\) (cf.~V, p.~71, exerc. 8).
\end{enumerate}

\begin{enumerate}
\def\labelenumi{\arabic{enumi})}
\setcounter{enumi}{12}
\item
  Let \(\mathbf{E}\) be a Banach space, and \(\mathbf{M}\) a closed vector subspace of \(\mathbf{E}\) . Show that if \(\mathbf{M}\) and \(\mathrm{E / M}\) are reflexive, then \(\mathbf{E}\) is reflexive.
\item
  Let \(\mathbf{A}\) be an infinite set.
\end{enumerate}

\begin{enumerate}
\def\labelenumi{\alph{enumi})}
\tightlist
\item
  Show that the strong dual of the Banach space \(E = \overline{\mathcal{K}(A)}\) (IV, p.~47, exerc. 1) can be identified with the Banach space \(\ell^{1}(A)\) , and that the strong dual of the Banach space \(\ell'(A)\) can be identified with the Banach space \(\mathcal{B}(A) = \ell^{\infty}(A)\) ; deduce that E is not reflexive and that \(E''/E\) is infinite dimensional (cf.~IV, p.~71, exerc. 18). If A = N, then E and \(E'\) are Banach spaces which satisfy the first axiom of countability, but \(E''\) does not (I, p.~25, exerc. 1).
\item
  Let \(B''\) be the unit ball in \(E'' = \ell^{\infty}(A)\) , and let \(B_{0}''\) be the convex set \(B'' + (B'' \cap E)\) . Show that \(B_{0}''\) is a bounded closed convex set in \(E''\) for the strong topology, with a non-empty interior, but does not have any extremal point. If p is the gauge of \(B_{0}''\) , show that \(E''\) endowed with the norm p is not isometric to any dual of a Banach space and that \(B_{0}''\) is not closed for the topology \(\sigma(E'', E')\) (although \(B''\) is compact for \(\sigma(E'', E')\) and \(B'' \cap E\) is strongly closed).
\end{enumerate}

¶ 15) The notations are those of exerc. 14.

\begin{enumerate}
\def\labelenumi{\alph{enumi})}
\tightlist
\item
  Let \((x_n')\) be a sequence in \(E'\) which converges to 0 for the topology \(\sigma(E', E'')\) ; show that, for every \(\varepsilon > 0\) , there exists a finite subset \(H\) of \(A\) such that \(\sum_{\alpha \notin H} |x_n'(\alpha)| \leqslant \varepsilon\) for every integer \(n\) .
\end{enumerate}

(Argue by reducto ad absurdum: if the property were not true, show that then there exists a number \(\delta > 0\) , an increasing sequence \((n_k)\) of integers, an increasing sequence \((\mathrm{H}_k)\) of finite subsets of A, such that \(\sum_{\alpha \in \mathrm{H}_{k-1}} |x_n'(\alpha)| \leqslant \frac{\delta}{8}\) for \(n \geqslant n_k\) , \(\sum_{\alpha \notin \mathrm{H}_k} |x_n'(\alpha)| \leqslant \frac{\delta}{8}\) for \(n \leqslant n_k\) and

\(\sum_{\alpha \in \mathrm{H}_k - \mathrm{H}_{k-1}} |x_n'(\alpha)| \geqslant \frac{\delta}{2}\) ; show that this implies a contradiction (the « gliding bump » method).)

Deduce that the sequence \((x_{n}^{\prime})\) converges to 0 for the strong topology, although the latter is strictly finer than the topology \(\sigma (\mathrm{E}',\mathrm{E}'')\) .

\begin{enumerate}
\def\labelenumi{\alph{enumi})}
\setcounter{enumi}{1}
\tightlist
\item
  Show that if \((x_{n}^{\prime})\) is a Cauchy sequence in \(E'\) for the topology \(\sigma(E', E'')\) , it converges to a point in \(E'\) for this topology; in other words, \(E'\) is semi-complete (III, p.~7) for \(\sigma(E', E'')\) . (Show that, for every \(\varepsilon > 0\) , there exists a finite subset \(H\) in \(A\) such that \(\sum_{\alpha \notin H} |x_{n}'(\alpha)| \leqslant \varepsilon\) for every integer n; argue by reducto ad absurdum, as in a), and use a).
\end{enumerate}

¶ 16) With the notations of exerc. 14, let E'' be the dual of \(E'' = \ell^{\infty}(A)\).

\begin{enumerate}
\def\labelenumi{\alph{enumi})}
\item
  Let \(e_{\alpha}\) (for \(\alpha \in A\) ) be the element of \(E''\) for which \(e_{\alpha}(\beta) = \delta_{\alpha \beta}\) (Kronecker's symbol). Let \((\mathbf{K}_n)\) be a sequence of finite subsets of \(A\) , two by two disjoint, and let \((x_n''')\) be a sequence of points in \(E'''\) . Show that there exists \((n_k)\) , a strictly increasing infinite sequence of integers \(>0\) such that, if we put \(B = \bigcup_k K_{n_k}\) , then the elements \(y_n' = (x_n'''(\alpha))_{\alpha \in B}\) belong to \(\ell^1(B)\) . (Let \(\delta\) be an arbitrary number \(>0\) , and let \((\mathrm{J}_m)_{m \in \mathbb{N}}\) be a partition of \(\mathbf{N}\) in finite sets. Arguing by reducto ad absurdum, show that if \(x''' \in E'''\) , then there exists an integer \(m\) such that \(|\langle x'', x''' \rangle| \leqslant \delta\) for all \(x'' \in E''\) for which \(\| x'' \| \leqslant 1\) and \(x''(\alpha) = 0\) except for indices \(\alpha\) belonging to the set \(\bigcup_{n \in J_m} K_n\) . Apply this result successively to \(x_1'', x_2'', ...\) in a suitable way.)
\item
  Deduce from \(a\) and from exerc. 15, a) that, if \((x_{n}^{\prime \prime \prime})\) is a sequence converging to 0 in \(\mathrm{E}'''\) for the topology \(\sigma (\mathrm{E}''',\mathrm{E}'')\) and if \(\tilde{x}_n^{\prime \prime \prime}\) is the restriction of \(x_{n}^{\prime \prime \prime}\) to the strongly closed subspace E of \(\mathrm{E}''\) , then \(\lim_{n\to \infty}\| \tilde{x}_n^{\prime \prime \prime}\| = 0\) in \(\mathrm{E}'\) .
\item
  Deduce from b) that, the strongly closed subspace E of \(E''\) does not have a topological complement for the strong topology. (Restricting to the case A = N, let \((e_{n}')\) be the sequence of continuous linear forms on E such that \(\langle x, e_{n}' \rangle = x(n)\) for all \(x \in E\) ; show that the sequence \((e_{n}')\) tends to 0 for \(\sigma(\mathrm{E}', \mathrm{E})\) , but that \(e_{n}'\) cannot be extended to a continuous linear form \(x_{n}'''\) on \(E''\) in such a way that the sequence \((x_{n}'''')\) tends to 0 for \(\sigma(\mathrm{E}'''', \mathrm{E}'')\) .)
\end{enumerate}

\begin{enumerate}
\def\labelenumi{\arabic{enumi})}
\setcounter{enumi}{16}
\tightlist
\item
  Let E be a non-reflexive Banach space, \(E'\) its strong dual, \(E''\) the strong dual of \(E'\) , \(E'''\) the strong dual of \(E''\) and \(E^{IV}\) the strong dual of \(E'''\) .
\end{enumerate}

\begin{enumerate}
\def\labelenumi{\alph{enumi})}
\item
  Show that, in \(E^{\prime\prime}\) , \(E^{\prime}\) and the subspace \(E^{\circ}\) orthogonal to E (when E is considered as a subspace of \(E^{\prime\prime}\) ) are topological complements, and that the projection from \(E^{\prime\prime}\) onto \(E^{\prime}\) for this decomposition is a continuous linear mapping of norm 1. Comparing with exerc. 16, c), deduce that the Banach space \(\mathcal{K}(A)\) is not isomorphic (as a topological vector space) to the strong dual of any Banach space.
\item
  Show that \(E^{IV}\) is the topological direct sum of \(E^{\circ}\) and \(E''\) , and also of \(E^{\prime\circ}\) and \(E^{\circ\circ}\) ; we have \(E'' \cap E^{\circ\circ} = E\) . Let v be the linear mapping from \(E^{IV}\) onto itself which is identity on \(E^{\circ}\) , and on \(E''\) , is the projection from \(E''\) onto \(E^{\circ\circ}\) parallel to \(E^{\circ}\) ; show that v is an isometry, but is not continuous for the topology \(\sigma(\mathrm{E}^{\mathrm{IV}}, \mathrm{E}''')\) .
\end{enumerate}

\begin{enumerate}
\def\labelenumi{\arabic{enumi})}
\setcounter{enumi}{17}
\item
  Show that a Banach space E whose strong dual \(E'\) satisfies the first axiom of countability, and which is semi-complete (III, p.~7) for the weakened topology \(\sigma(\mathrm{E}, \mathrm{E}')\) , is reflexive (compare with IV, p.~54, exerc. 15, b)).
\item
  Let \(\mathrm{E}\) be a Banach space, \(\mathrm{E}'\) its strong dual, \(\mathrm{G}'\) a strongly closed subspace of \(\mathrm{E}'\) satisfying the first axiom of countability for the strong topology. Show that there exists a countable subset of \(\mathrm{E}\) such that, if \(\mathrm{F}\) is the closed vector subspace of \(\mathrm{E}\) generated by this subset, then \(\mathrm{G}'\) is isometric to a strongly closed subspace of the strong dual \(\mathrm{F}'\) of \(\mathrm{F}\) . (Suppose that the sequence \((x_n')\) is strongly dense in \(\mathrm{G}'\) ; for every \(n\) , let \(x_n \in \mathrm{E}\) such that \(\| x_n \| \leqslant 1\) and \(\langle x_n, x_n' \rangle = \left(1 - \frac{1}{n}\right) \| x_n' \|\) ; show that the strongly closed subspace \(F\) of \(E\) generated by the
\end{enumerate}

\[
x _ {n}
\]

is the required space.)

¶ 20) Let E be a Banach space, E' its dual, and B the unit ball in E. In order that every point of B have a countable fundamental system of neighbourhoods for the topology induced by the weakened topology σ(E, E') on B, it is necessary and sufficient that E' satisfies the first axiom of countability for the strong topology. (To show that the condition is necessary, observe that if every point in B has a countable fundamental system of neighbourhoods for the weakened topology, then this is also true for the closure B°° of B in E'' for the topology σ(E'', E'). Therefore there exists a sequence (a'\_n) in E' such that every neighbourhood of 0 in B°° for σ(E'', E') contains the intersection of B°° and of a finite number of polars \{a'\_n\}°; consider the strongly closed subspace W' of E' generated by the a'\_n, and the orthogonal \(W'^{\circ}\) of W' in E''.)

\begin{enumerate}
\def\labelenumi{\arabic{enumi})}
\setcounter{enumi}{20}
\tightlist
\item
  Let E be a Banach space, \(E'\) its dual, B the unit ball in E and \(B_{r}'\) the closed ball with centre 0 and radius r in \(E'\) .
\end{enumerate}

\begin{enumerate}
\def\labelenumi{\alph{enumi})}
\item
  Let \(M_{1}^{\prime}\) , \(M_{2}^{\prime}\) be two vector subspaces of \(E^{\prime}\) , which are everywhere dense for the weak topology \(\sigma(\mathrm{E}^{\prime}, \mathrm{E})\) . In order that the topologies induced by \(\sigma(\mathrm{E}, \mathrm{M}_{1}^{\prime})\) and \(\sigma(\mathrm{E}, \mathrm{M}_{2}^{\prime})\) on B coincide, it is necessary and sufficient that the strong closures of \(M_{1}^{\prime}\) and \(M_{2}^{\prime}\) in \(E^{\prime}\) are identical.
\item
  Let \(M'\) be a vector subspace of \(E'\) which is everywhere dense for \(\sigma(E', E)\) ; let \(\mathbf{M}'^{(1)}\) denote the vector subspace generated by the closure of \(M' \cap B_{1}'\) in \(E'\) for the topology \(\sigma(E', E)\) . In order that \(\mathbf{M}'^{(1)} = E'\) , it is necessary and sufficient that the weak closure of \(M' \cap B_{1}'\) contains a ball \(B_{r}'\) with r \textgreater{} 0 (use the fact that \(E'\) is barrelled for the strong topology).
\item
  Let \(r\) be the supremum of the numbers \(t\) such that the weak closure of \(\mathbf{M}' \cap \mathbf{B}_1'\) contains a ball \(\mathbf{B}_t'\) ; the number \(r\) is said to be the characteristic of \(\mathbf{M}'\) . Show that \(r\) is the infimum of the numbers \(\sup_{x' \in \mathbf{M}' \cap \mathbf{B}_1'} |\langle x, x' \rangle| / \| x\|\) where \(x\) ranges over the set of all points \(\neq 0\) in \(E\) (use the Hahn-Banach theorem).
\item
  Show that \(1 / r\) is the supremum of \(\| x\|\) as \(x\) ranges over the closure of B in E for the topology \(\sigma (\mathrm{E},\mathbf{M}^{\prime})\) (use \(c\) ) and the Hahn-Banach theorem).
\item
  Let \(M^{\prime\circ}\) be the orthogonal of \(M^{\prime}\) in \(E^{\prime\prime}\) ; show that \(r = \inf(\|x + z^{\prime\prime}\|/\|x\|)\) where \(z^{\prime\prime}\) ranges over \(M^{\prime\circ}\) and x ranges over the set of non-zero points in E (use c) and the Hahn-Banach th). Deduce that in order that \(\mathbf{M}^{\prime(1)} = \mathbf{E}^{\prime}\) , it is necessary and sufficient that \(E + M^{\prime\circ}\) be strongly closed in \(E^{\prime\prime}\) (use Banach's th., I, p.~17).
\end{enumerate}

\(f\) ) Let \(\mathbf{A} = \mathbf{N}\times \mathbf{N}\) and \(\mathrm{E} = \overline{\mathcal{K}(\mathbf{A})}\) (cf.~IV, p.~54, exerc. 14). In the space \(\mathrm{E}'' = \ell^{\infty}(\mathrm{A})\) , let \(\mathrm{P}\) be the vector subspace consisting of all points \(x = (x_{ij})\) such that \(x_{ij} = x_{0j} / (j + 1)\) for all \(i\geqslant 0\) . Show that \(\mathbf{P} = \mathbf{M}'^{\circ}\) where \(\mathbf{M}'\) is a vector subspace of \(\mathrm{E}'\) that is everywhere dense (for \(\sigma (\mathrm{E}',\mathrm{E}))\) , but that \(\mathrm{E} + \mathrm{M}'^{\circ} = \mathrm{E} + \mathrm{P}\) is not strongly closed in \(\mathrm{E}''\) ; deduce that the characteristic of \(\mathbf{M}'\) is 0.

\begin{enumerate}
\def\labelenumi{\arabic{enumi})}
\setcounter{enumi}{21}
\tightlist
\item
  Let E be a Banach space, \(E'\) its dual and \(M'\) a strongly closed subspace in \(E'\) which is everywhere dense for the weak topology on \(E'\) . We say that M is irreducible if there exists no vector subspace \(N' \neq M'\) of \(M'\) which is strongly closed and weakly everywhere dense in \(E'\) . a) Show that \(M'\) is irreducible if and only if the orthogonal \(M'^{\circ}\) of \(M'\) in \(E''\) is the topological complement of E (for the strong topology of \(E''\) ). Deduce that then \(\mathbf{M}^{(1)} = \mathbf{E}'\) (exerc. 21) and that E is isomorphic to the strong dual of the space \(M'\) endowed with the topology induced by the strong topology of \(E'\) .
\end{enumerate}

\begin{enumerate}
\def\labelenumi{\alph{enumi})}
\setcounter{enumi}{1}
\item
  Show that \(\mathbf{M}'\) is irreducible if and only if the unit ball in \(\mathbf{E}\) is relatively compact for the topology \(\sigma (\mathrm{E},\mathbf{M}')\) (use exerc. 21, a)).
\item
  For E to be isomorphic to a strong dual of a Banach space (for the topological vector space structure), it is necessary and sufficient that there exist an irreducible subspace in \(E'\) . Deduce a new proof of the fact that the Banach space \(\overline{\mathcal{K}(\mathbf{N})}\) is not isomorphic to a strong dual of a Banach space (cf.~IV, p.~55, exerc. 16, c)).
\end{enumerate}

\begin{enumerate}
\def\labelenumi{\arabic{enumi})}
\setcounter{enumi}{22}
\tightlist
\item
  With the same notations as in exerc. 22, assume that \(\mathbf{M}'\) is irreducible.
\end{enumerate}

\begin{enumerate}
\def\labelenumi{\alph{enumi})}
\item
  In order that the canonical mapping from E into \(E^{\prime\prime}/M'^{\circ}\) , which is the restriction of the canonical mapping \(E^{\prime\prime} \rightarrow E^{\prime\prime}/M'^{\circ}\) be a Banach space isometry, it is necessary and sufficient that the characteristic (IV, p.~55, exerc. 21) of \(M^{\prime}\) is equal to 1. Then, \(M^{\prime}\) endowed with the norm induced by that of \(E^{\prime}\) is said to be the predual of E and E can be canonically identified (with its norm) with the dual of the Banach space \(M^{\prime}\) .
\item
  For every vector subspace F of E which is closed for \(\sigma (\mathrm{E},\mathbf{M}^{\prime})\) , show that the canonical image of \(\mathbf{M}'\) in the dual \(\mathrm{F}'\) of F, identified with \(\mathrm{E}^{\prime} / \mathrm{F}^{\circ}\) , is a predual of F.
\end{enumerate}

\begin{enumerate}
\def\labelenumi{\arabic{enumi})}
\setcounter{enumi}{23}
\tightlist
\item
  \begin{enumerate}
  \def\labelenumii{\alph{enumii})}
  \tightlist
  \item
    Let \((a_{k})_{1\leqslant k\leqslant n}\) be a finite sequence of points in a normed space \(E\) , and let \((\lambda_k)_{1\leqslant k\leqslant n}\) be a finite sequence of numbers \(>0\) such that \(\sum_{k = 1}^{n}\lambda_{k} < 1\) ; put \(\mu_{k} = 1 - \sum_{j = 1}^{k - 1}\lambda_{j}\) for every \(k\) ; then
  \end{enumerate}
\end{enumerate}

\[
\left\| \sum_ {k = 1} ^ {n} \lambda_ {k} a _ {k} \right\| \leqslant \frac {\lambda_ {n}}{\mu_ {n - 1}} \left\| \mu_ {n - 1} a _ {n} + \sum_ {k = 1} ^ {n - 1} \lambda_ {k} a _ {k} \right\| + \frac {\mu_ {n}}{\mu_ {n - 1}} \left\| \sum_ {k = 1} ^ {n - 1} \lambda_ {k} a _ {k} \right\|.
\]

\begin{enumerate}
\def\labelenumi{\alph{enumi})}
\setcounter{enumi}{1}
\tightlist
\item
  Let \((a_{n})\) be an infinite sequence of points in the unit ball of E, and let \((\lambda_{n})\) be an infinite sequence of numbers \(>0\) such that \(\sum_{n} \lambda_{n} = 1\) . For every \(n > 0\) , put \(\mu_{n} = 1 - \sum_{k=1}^{n-1} \lambda_{k}\) and \(b_{n} = \sum_{k=1}^{n-1} \lambda_{k} a_{k} + \mu_{n} a_{n}\) ; show that, for all \(n \geqslant 1\) , we have
\end{enumerate}

\[
\left\| \sum_ {k = 1} ^ {n} \lambda_ {k} a _ {k} \right\| \leqslant \mu_ {n} \sum_ {k = 1} ^ {n} \frac {\lambda_ {k} \| b _ {k} \|}{\mu_ {k - 1} \mu_ {k}}.
\]

(Apply a) by induction.)

\begin{enumerate}
\def\labelenumi{\alph{enumi})}
\setcounter{enumi}{2}
\tightlist
\item
  Let \((\mathrm{C}_n)\) be a decreasing sequence of convex sets in E, contained in the unit ball, and suppose that \(d(0,\mathrm{C}_1)\geqslant \theta >0\) (hence, a fortiori \(d(0,\mathrm{C}_n)\geqslant \theta\) for all \(n\) ). Let \((\lambda_{n})\) be a sequence of numbers \(>0\) such that \(\sum_{n}\lambda_{n} = 1\) . Show that there exists a number \(\alpha\) such that \(\theta \leqslant \alpha \leqslant 1\) and a sequence \((x_{n})\) of points of E such that \(x_{n}\in \mathbf{C}_{n}\) for all \(n\) , \(\| \sum_{n}\lambda_{n}x_{n}\| = \alpha\) and, for all \(n\)
\end{enumerate}

\[
\left\| \sum_ {k = 1} ^ {n} \lambda_ {k} x _ {k} \right\| \leqslant \alpha (1 - \theta \sum_ {j = n + 1} ^ {\infty} \lambda_ {j}).
\]

(Take \(x_{1}\) such that \(\| x_{1}\|\) is arbitrarily close to \(d(0, \mathrm{C}_{1})\) , then, by induction, take \(x_{n}\) such that \(\left\| \sum_{k=1}^{n-1} \lambda_{k} x_{k} + \left(\sum_{j=n}^{\infty} \lambda_{j}\right) x_{n}\right\|\) is arbitrarily close to the infimum of the numbers \(\left\| \sum_{k=1}^{n-1} \lambda_{k} x_{k} + \left(\sum_{j=n}^{\infty} \lambda_{j}\right) y\right\|\) where \(y\) ranges over \(\mathrm{C}_n\) . Then use \(b)\) .)

¶ 25) Let E be a Banach space satisfying the first axiom of countability. Show that the following properties are equivalent :

α) E is not reflexive.

\(\beta)\) For every number \(\theta\) such that \(0 < \theta < 1\) , there exists a sequence \((x_n')\) in \(\mathbf{E}'\) such that, \(\| x_n'\| \leqslant 1\) for all \(n\) , that the sequence \((x_n')\) converges to 0 for \(\sigma (\mathrm{E}',\mathrm{E})\) and that the distance of 0 from the convex set generated by the \(x_{n}^{\prime}\) is \(\geqslant \theta\) .

\(\gamma)\) For every number \(\theta\) such that \(0 < \theta < 1\) and every sequence \(\lambda_{n}\) of numbers \(>0\) such that \(\sum_{n} \lambda_{n} = 1\) , there exists a number \(\alpha\) such that \(\theta \leqslant \alpha \leqslant 1\) and a sequence \((y_{n}')\) of points of

E' such that \(\| y_n'\| \leqslant 1\) for all \(n\) , \(\left\| \sum_{n}\lambda_ny_n'\right\| = \alpha\) and \(\left\| \sum_{k=1}^{n}\lambda_ky_k'\right\| \leqslant \alpha (1 - \theta \sum_{j=n+1}^{\infty}\lambda_j)\) for all \(n\) .

\(\delta)\) There exists \(z' \in \mathrm{E}'\) such that for no \(x \in \mathrm{E}\) do we have \(\left|\langle x, z' \rangle\right| = \| x \| \cdot \| z'\|\) (Theorem of James-Klee).

(To see that \(\alpha\) ) implies \(\beta\) ), observe that there exists \(z'' \in \mathrm{E}''\) such that \(\| z'' \| < 1\) and \(d(z'', \mathrm{E}) > \theta\) . If \((x_n)\) is an everywhere dense sequence in \(\mathrm{E}\) , find the sequence \((x_n')\) in \(\mathrm{E}'\) such that \(\| x_n' \| < 1\) and such that \(\langle x_k, x_n' \rangle = 0\) for \(k \leqslant n\) and \(\langle x_n', z'' \rangle = \theta\) . To see that \(\beta\) ) implies \(\gamma\) ), use exerc. 24. For \(\gamma\) ) implies \(\delta\) ), show that for all \(x \in \mathrm{E}\) we have \(|\sum_{n} \lambda_n \langle x, y_n' \rangle| < \alpha\) , with the notations of \(\gamma\) ).

\begin{enumerate}
\def\labelenumi{\arabic{enumi})}
\setcounter{enumi}{25}
\tightlist
\item
  A locally convex space E is said to have the property (GDF) if every linear mapping u from E into a Banach space F which satisfies the following property, is continuous: in the product space \(E \times F\) , every limit of a convergent sequence of points in the graph \(\Gamma\) of u again belongs to \(\Gamma\) . Every Fréchet space has the property (GDF) (I, p.~19, cor. 5); this is also true of every inductive limit of a family of Fréchet spaces (II, p.~34, prop. 10). Show that every Hausdorff locally convex space with the (GDF) property is barrelled. (Let V be a barrel in E, q its gauge and H the Hausdorff space associated with E endowed with this semi-norm; show that the canonical mapping \(\pi\) from E into the completion \(\hat{H}\) is continuous by using the property (GDF) and the fact that every linear form \(x' \in V^{0}\) can be extended uniquely to a continuous linear form on \(\hat{H}\) , the set of these forms being the unit ball in the dual of \(\hat{H}\) .)
\end{enumerate}

\exercisegroup

¶ 1) Let E be a locally convex metrizable space, and \(E_{b}^{\prime}\) its strong dual. If \(E_{b}^{\prime}\) is metrizable, prove that the topology of E can be defined by a single norm (use III, p.~37, exerc. 2 and p.~38, exerc. 5 and also the fact that E is bornological).

\begin{enumerate}
\def\labelenumi{\arabic{enumi})}
\setcounter{enumi}{1}
\tightlist
\item
  An infra-barrelled space is semi-barrelled. A locally convex space is said to be a (DF) space if it is semi-barrelled and if the canonical bornology (III, p.~3, def. 5) has a countable base. Every normed space and every strict inductive limit of a sequence of normed spaces (II, p.~33) is a (DF) space. Every strong dual of a Fréchet space is a (DF) space.
\end{enumerate}

\begin{enumerate}
\def\labelenumi{\alph{enumi})}
\item
  The strong dual of a (DF) space is a Fréchet space.
\item
  Let \(\mathrm{E}\) be a (DF) space and let \((\mathbf{A}_n)\) be an increasing sequence of bounded, convex, balanced and closed subsets of \(\mathrm{E}\) such that every bounded subset of \(\mathrm{E}\) is absorbed by one of the \(\mathbf{A}_n\) .
\end{enumerate}

Let U be the union of the \(A_{n}\) ; show that the closure \(\overline{U}\) of U in E is precisely the set of all \(x \in E\) such that \(\lambda x \in U\) for \(0 \leqslant \lambda < 1\) . (If \(x \notin \lambda U\) for some \(\lambda > 1\) , then for every n, there exists a linear form \(x_{n}' \in E'\) such that \(x_{n}' \in A_{n}^{\circ}\) and \(\langle x, x_{n}' \rangle = \lambda\) , and the sequence \((x_{n}')\) is equicontinuous, hence has a weak limit point.)

\begin{enumerate}
\def\labelenumi{\alph{enumi})}
\setcounter{enumi}{2}
\tightlist
\item
  Show that if a (DF) space is barrelled, it is also bornological (cf.~III, p.~44, exerc. 13, b)). Give examples of (DF) spaces which are not ultrabornological, but are bornological and barrelled (III, p.~46, exerc. 22) and also of (DF) spaces which are not barrelled but are bornological.
\end{enumerate}

\begin{enumerate}
\def\labelenumi{\arabic{enumi})}
\setcounter{enumi}{2}
\tightlist
\item
  Let \(E\) be a locally convex metrizable space, and \(E_b'\) its strong dual.
\end{enumerate}

\begin{enumerate}
\def\labelenumi{\alph{enumi})}
\item
  Show that every convex balanced subset \(V'\) of \(E_{b}'\) which absorbs the strongly bounded subsets of \(E_{b}'\) contains a barrel (for the strong topology) which absorbs the strongly bounded subsets of \(E_{b}'\) . (Let \((\mathbf{K}_{n}')\) be a countable base of the canonical bornology of \(E_{b}'\) and let \(\lambda_{n}\) be such that \(\lambda_{n}K_{n}' \subset \frac{1}{2}V'\) ; apply exerc. 2, b) to the sequence \(A_{n}'\) , where \(A_{n}'\) is the convex envelope of the union of the \(\lambda_{j}K_{j}'\) for \(j \leqslant n\) .)
\item
  Deduce from \(a\) ) that the following properties are equivalent :
\end{enumerate}

α) E is distinguished (IV, p.~52, exerc. 4).

β) \(E_{b}^{\prime}\) is infrabarrelled (III, p.~44, exerc. 7).

\(\gamma)\) \(\mathrm{E}_b^{\prime}\) is bornological.

δ) \(E_{b}^{\prime}\) is barrelled.

ε) \(E_{b}^{\prime}\) is ultrabornological (III, p.~45, exerc. 19).

\begin{enumerate}
\def\labelenumi{\alph{enumi})}
\setcounter{enumi}{2}
\tightlist
\item
  Show that if \(E_{b}^{\prime}\) is reflexive, then \(\hat{E} = E''\) (which is obviously reflexive) (cf.~IV, p.~52, exerc. 4 and p.~53, exerc. 11).
\end{enumerate}

\begin{enumerate}
\def\labelenumi{\arabic{enumi})}
\setcounter{enumi}{3}
\tightlist
\item
  Let E be a locally convex Hausdorff space, \(E'\) its dual. If M is a closed vector subspace of E which is metrizable and distinguished (IV, p.~52, exerc. 4), then the strong topology \(\beta(\mathrm{E}'/\mathrm{M}^{\circ}, \mathrm{M})\) is the quotient topology by \(M^{\circ}\) of the strong topology \(\beta(\mathrm{E}', \mathrm{E})\) (use exerc. 3, b) and IV, p.~51, exerc. 22, b)).
\end{enumerate}

¶ 5) For every integer n \textgreater{} 0, let \(a^{(n)}\) be the double sequence \((a_{pq}^{(n)})\) ( \(p \in N\) , \(q \in N\) ) such that \(a_{pq}^{(n)} = q\) if \(p \leqslant n\) and \(a_{pq}^{(n)} = 1\) if p \textgreater{} n.~Let E be the vector space of all double sequences \(x = (x_{pq})_{(p,q) \in \mathbf{N} \times \mathbf{N}}\) of real numbers such that, for every integer n \textgreater{} 0, the number \(r_n(x) = \sum_{p,q} a_{pq}^{(n)} |x_{pq}|\)

is finite. If E is assigned the topology defined by the semi-norms \(r_n\) , then E is a Fréchet space satisfying the first axiom of countability (IV, p.~47, exerc. 1, c)); the dual E' of E can be identified with the space of all double sequences \(x' = (x_{pq}')\) of real numbers such that for at least one index \(n\) , there exists \(k_n > 0\) such that \(|x_{pq}'| \leqslant k_n a_{pq}^{(n)}\) for every pair \((p, q)\) ; and \(\langle x, x' \rangle = \sum_{p, q} x_{pq} x_{pq}'\) .

(IV, p.~47, exerc. 1, c)).

For every integer \(p_0 > 0\) and every sequence \((m_p)\) of integers \(> 0\) , let \(\mathrm{J}(p_0; (m_p))\) be the set of pairs of integers \(p > 0\) , \(q > 0\) such that \(p \geqslant p_0\) and \(q \geqslant m_p\) ; let \(\mathfrak{V}\) be the filter base on \(\mathbf{N} \times \mathbf{N}\) consisting of the sets \(\mathrm{J}(p_0; (m_p))\) and let \(\mathfrak{F}\) be an ultrafilter which is finer than the filter with base \(\mathfrak{V}\) .

\begin{enumerate}
\def\labelenumi{\alph{enumi})}
\item
  Show that for all \(x' = (x_{pq}') \in \mathrm{E}'\) , the double sequence \((x_{pq}')\) has a limit \(u(x')\) with respect to the ultrafilter \(\mathfrak{F}\) ; if \(\mathrm{V}_n\) is a neighbourhood of 0 in \(\mathrm{E}\) defined by \(r_n(x) \leqslant 1\) , then \(|u(x')| \leqslant 1\) for all \(x' \in \mathrm{V}_n^\circ\) .
\item
  Let \(U'\) be a neighbourhood of 0 in \(E'\) , for the strong topology, which is convex, balanced and weakly closed, and for every n, let \(\alpha_{n} > 0\) be such that \(\alpha_{n} V_{n}^{\circ} \subset U'\) . For every integer p \textgreater{} 0, let \(m_{p}\) be an integer such that \(2^{p+1} \leqslant \alpha_{p} m_{p}\) , and let \(x' = (x'_{pq})\) be the double sequence with \(x'_{pq} = 0\) for q \textless{} \(m_{p}\) , \(x'_{pq} = 2\) for \(q \geqslant m_{p}\) . Show that \(x' \in U'\) but that \(u(x') = 2\) ; deduce that u is not strongly continuous in \(E'\) , while being bounded on every bounded subset of \(E'\) . Conclude (IV, p.~58, exerc. 3) that E is not distinguished, and consequently that the strong dual \(E'_{b}\) is a non infra-barrelled (DF) space.
\item
  Using b) construct an example of a closed subspace M of a Fréchet space F such that the strong topology \(\beta(\mathrm{F}^{\prime}/\mathrm{M}^{\circ}, \mathrm{M})\) is distinct from the quotient topology by \(M^{\circ}\) of the strong topology \(\beta(\mathrm{F}^{\prime}, \mathrm{F})\) (embed E in a countable product of Banach spaces).
\end{enumerate}

¶ 6) a) Let E be a (DF) space (IV, p.~57, exerc. 2), and let U be a convex set such that for every bounded, convex and balanced subset A of E, U ∩ A is a neighbourhood of 0 for the topology induced on A by that of E. Show that U is a neighbourhood of 0 in A. (Let (A \(_{n}\) ) be a countable base for the canonical bornology of E (III, p.~3, def. 5). Show that, by induction we can define a sequence ( \(\lambda_{n}\) ) of numbers \textgreater{} 0 and a sequence (V \(_{n}\) ) of closed, convex and balanced

neighbourhoods of 0 in E such that \(\lambda_{n}\mathrm{A}_{n} \subset \frac{1}{3}\mathrm{U}\) , \(\lambda_{n}\mathrm{A}_{n} \subset \bigcap_{j=1}^{\infty} \mathrm{V}_{j}\) , \(\mathrm{V}_{n} \cap \mathrm{A}_{n} \subset \mathrm{U}\) for every \(n\) .

First show that if \(\lambda_{j}\) and \(\mathrm{V}_j\) have been constructed for \(j\leqslant n\) , then we can find \(\lambda_{n + 1}\) such that \(\lambda_{n + 1}\mathrm{A}_{n + 1}\subset \frac{1}{3}\mathrm{U}\) and \(\lambda_{n + 1}\mathrm{A}_{n + 1}\subset \mathrm{V}_j\) for \(j\leqslant n\) . Next prove that we can find \(\mathrm{V}_{n + 1}\) such that \(\lambda_j\mathrm{A}_j\subset \mathrm{V}_{n + 1}\) for \(j\leqslant n + 1\) and \(\mathrm{V}_{n + 1}\cap \mathrm{A}_{n + 1}\subset \mathrm{U}\) ; for this, letting A denote the convex envelope of the \(\lambda_{j}\mathrm{A}_{j}\) for \(j\leqslant n + 1\) , show that we can take \(\mathrm{V}_{n + 1} = \overline{\mathrm{A} + \mathrm{V}}\) for a suitable convex balanced neighbourhood V of 0; we remark that for this it is enough to show that, if \(\mathbf{B} = \mathbf{A}_{n + 1}\cap \mathbb{C}\) U, then 0 is not in the closure of the set \(\mathbf{B} + 2\mathbf{A}\) .

\begin{enumerate}
\def\labelenumi{\alph{enumi})}
\setcounter{enumi}{1}
\tightlist
\item
  Deduce from a) that if u is a linear mapping from E into a locally convex space F, such that the restriction of u to every bounded subset of E is continuous, then u is continuous (cf.~IV, p.~50, exerc. 15).
\end{enumerate}

¶ 7) a) Let E be a (DF) space, U a convex, balanced and closed set in E, which absorbs the bounded subsets of E, and let \((x_n)\) be a sequence of points of \(\complement\) U. Show that there exists a neighbourhood V of 0 in E which does not contain any of the \(x_n\) . (Let \((A_n)\) be a countable base for the canonical bornology of E. Show that, by induction we can define a sequence \((\lambda_n)\) of numbers \textgreater{} 0 and a sequence \((V_n)\) of convex, balanced and closed neighbourhoods of

0 such that \(\lambda_{n}\mathbf{A}_{n}\subset \bigcap_{j = 1}\mathbf{V}_{j},\lambda_{n}\mathbf{A}_{n}\subset \mathrm{U}\) and \(x_{n}\in \mathbb{C}\mathbf{V}_{n}\) for all \(n\) . For this, if the \(\lambda_{j}\) and \(\mathrm{V}_j\) have

been constructed for \(j \leqslant n\) , take \(\lambda_{n+1}\) such that \(\lambda_{n+1}\mathrm{A}_{n+1} \subset \mathrm{U}\) and \(\lambda_{n+1}\mathrm{A}_{n+1} \subset \mathrm{V}_j\) for all \(j \leqslant n\) , then take \(\mathrm{V}_{n+1}\) containing the closure of the convex envelope of the union of the \(\lambda_j\mathrm{A}_j\) for \(j \leqslant n + 1\) .

\begin{enumerate}
\def\labelenumi{\alph{enumi})}
\setcounter{enumi}{1}
\item
  Deduce from a) that if M is a subset of E containing an everywhere dense countable set, then the topology induced on M by the strong topology of the bidual \(E''\) of E is identical with the topology induced by that of E. In particular, the convergent sequences in E are the same for the topology of E and for the topology induced by the strong topology of \(E''\) ; for every metrizable subset M of E, the topology induced on M by the topology of E is identical with the topology induced by the strong topology of \(E''\) .
\item
  Deduce from \(a\) ) that if there exists a countable everywhere dense set in \(E\) , then \(E\) is infra-barrelled.
\item
  Deduce from b) and from exerc. 6 that if every bounded subset of E is metrizable for the topology induced by that of E, then E is infrabarrelled.
\end{enumerate}

\begin{enumerate}
\def\labelenumi{\arabic{enumi})}
\setcounter{enumi}{7}
\tightlist
\item
  Let E be a Fréchet space, \(E_{b}^{\prime}\) its strong dual. Suppose that there exists an everywhere dense sequence \((x_{n}^{\prime})\) in \(E_{b}^{\prime}\) . Show that E satisfies the first axiom of countability. (Let \((\mathbf{K}_{n}^{\prime})\) be a countable base for the canonical bornology of \(E_{b}^{\prime}\) consisting of closed convex balanced sets. For every system \(\alpha\) consisting of a point \(x_{n}^{\prime}\) , an arbitrary finite number of rational numbers \(\lambda_{k} > 0 (1 \leqslant k \leqslant m)\) and m indices \(n_{k}\) such that \(x_{n}^{\prime} \notin 2 \sum_{k=1}^{m} \lambda_{k} K_{n_{k}}^{\prime} = 2 H_{\alpha}^{\prime}\) , let \(x_{\alpha} \in E\) be such that
\end{enumerate}

the hyperplane with equation \(\langle x_{\alpha}, y' \rangle = 1\) strictly separates the two weakly compact sets \(\mathrm{H}_{\alpha}'\) and \(x_{n}' + \mathrm{H}_{\alpha}'\) . Show that for every \(x' \neq 0\) in \(\mathrm{E}'\) , there exists a system \(\alpha\) such that \(\langle x_{\alpha}, x' \rangle \neq 0\) . For this, consider a neighbourhood \(\mathrm{V}'\) of 0 in \(\mathrm{E}_b'\) such that \(\mathrm{V}' \cap (x' + \mathrm{V}') = \emptyset\) , then for each integer \(m\) , take a rational number \(\lambda_m > 0\) such that \(\lambda_m \mathrm{K}_m' \subset \mathrm{V}'\) ; use the fact that the union \(\mathrm{U}' \subset \mathrm{V}'\) of the \(\lambda_m \mathrm{K}_m'\) is a neighbourhood of 0 (exerc. 7, and IV, p.~58, exerc. 3, b)) and that there exists \(n\) such that \(x_n' \in x' + \mathrm{U}'\) .

¶ 9) a) Let E be a Hausdorff semi-barrelled space, M a closed vector subspace of E and E' the dual of E. Show that E/M is semi-barrelled and that the strong topology β(M°, E/M) is identical with the topology induced on M° by the strong topology β(E', E). (Note that it is enough to prove that a sequence (x') in M° which converges to 0 for β(E', E) is bounded for β(M°, E/M).) Deduce that if E is a (DF) space, then so is E/M (cf.~IV, p.~63, exerc. 8).

\begin{enumerate}
\def\labelenumi{\alph{enumi})}
\setcounter{enumi}{1}
\item
  Let E be a locally convex Hausdorff space, M a (not-necessarily closed) vector subspace of E. Show that if M is a semi-barrelled space, then the strong topology \(\beta(\mathrm{E}^{\prime}/\mathrm{M}^{\circ}, \mathrm{M})\) is identical with the quotient topology by \(M^{\circ}\) of the strong topology \(\beta(\mathrm{E}^{\prime}, \mathrm{E})\) . (Argue as in a).)
\item
  Show that a Hausdorff and quasi-complete semi-barrelled space E is complete (use b) applied to E and \(\hat{E}\) ). In particular, a semi-barrelled, semi-reflexive space is complete (cf.~IV, p.~52, exerc. 6).
\item
  Show that the completion of a semi-barrelled (resp. (DF)) Hausdorff space is semi-barrelled (resp. a (DF) space).
\item
  Let \((\mathrm{E}_n)\) be a sequence of semi-barrelled (resp. (DF)) spaces, E a vector space, and for each \(n\) , let \(f_n\) be a linear mapping from \(\mathrm{E}_n\) into E. Suppose that \(\dot{\mathrm{E}}\) is the union of the \(f_n(\mathrm{E}_n)\) ; show that, E is semi-barrelled (resp. a (DF) space) for the finest locally convex topology for which all the \(f_n\) are continuous (first examine the case where E is the topological direct sum of the \(\mathrm{E}_n\) ). If the \(\mathrm{E}_n\) are semi-reflexive (resp. reflexive) and if E is Hausdorff, then E is semi-reflexive (resp. reflexive).
\end{enumerate}

\begin{enumerate}
\def\labelenumi{\arabic{enumi})}
\setcounter{enumi}{9}
\tightlist
\item
  Let E be a Fréchet space satisfying the first axiom of countability. Show that if in the dual \(E'\) of E, every sequence which converges for the weak topology \(\sigma(E', E)\) also converges for the strong topology \(\beta(E', E)\) , then E is a Montel space. (Show that every bounded subset of \(E'\) is relatively compact for the strong topology; use GT, II, § 4, exerc. 6; then use IV, p.~53, exerc. 11, b).)
\end{enumerate}

¶ 11) Let \((c_{mn})\) be a double sequence of numbers \textgreater{} 0 such that \(c_{m,n} \leqslant c_{m+1,n}\) and let E be the space of all sequences \(x = (x_n)\) of real numbers such that \(p_m(x) = \sum_n c_{mn} |x_n| < +\infty\) for

every integer \(m\) . We endow \(E\) with the topology defined by the semi-norms \(p_m\) and for this topology \(E\) is a Fréchet space satisfying the first axiom of countability; the dual \(E'\) of \(E\) can be identified with the space of all sequences \(x' = (x_n')\) such that \(\sup_{n} c_{mn}^{-1} |x_n'| < +\infty\) for at

least one \(m\) , the canonical bilinear form \(\langle x, x' \rangle\) being identified with \(\sum_{n} x_n x_n'\) (IV, p.~47,

exerc. 1, c)). Suppose that there does not exist any subsequence \((n_k)\) for which there is a sequence \((a_m)\) of numbers \(\geqslant 0\) and an index \(m_0\) such that \(c_{m,n_k} \leqslant a_m c_{m_0,n_k}\) for all \(m \geqslant m_0\) and for all \(k\) . Under these conditions, every weakly convergent sequence in \(\mathbf{E}'\) is strongly convergent and consequently (IV, p.~60, exerc. 10) E is a Montel space. (Argue by reductio ad absurdum; if necessary make a transformation of the form \((x_n) \mapsto (a_n x_n)\) to reduce to the case where \(c_{m_0 n} = 1\) for all \(n\) and some \(m_0\) , and where there exists a sequence \((x'^{(p)})_{p \geqslant 0}\) which converges weakly to 0 in \(\mathbf{E}'\) , and is such that \(|x_n'^{(p)}| \leqslant 1\) for every pair \((p, n)\) , and also that there exists a bounded set B in E, defined by \(p_m(x) \leqslant b_m\) for all \(m \geqslant 0\) and such that \(\sup_{x \in \mathbf{B}} |\langle x, x'^{(p)} \rangle| \geqslant 2\delta > 0\) for

every integer \(p\) . Under these hypothesis, prove that there exists a strictly increasing sequence \((r_q)\) of integers, and a sequence \((x^{(q)})\) of points of \(\mathbf{B}\) such that

\[
\sum_ {k = r _ {q} + 1} ^ {r _ {q+1}} \left| x _ {k} ^ {(q)} \right| > \delta
\]

for each index \(q\) . Then show, by reductio ad absurdum, that for every \(q\) , there exists at least one index \(s_q\) such that \(r_q < s_q \leqslant r_{q+1}\) and that for every integer \(m\) , we have \(c_{m,s_q} \leqslant b_m 2^{m+2}/\delta\) , which contradicts the hypothesis.)

\begin{enumerate}
\def\labelenumi{\arabic{enumi})}
\setcounter{enumi}{11}
\tightlist
\item
  \begin{enumerate}
  \def\labelenumii{\alph{enumii})}
  \tightlist
  \item
    Let F be a Hausdorff (DF) space, and \(F_{b}^{\prime}\) its strong dual. Show that if \(F_{b}^{\prime}\) is reflexive, then the completion \(\hat{F}\) of F is reflexive and is equal to the bidual \(F^{\prime\prime}\) of F (cf.~IV, p.~52, exerc. 4 and p.~53, exerc. 11).
  \end{enumerate}
\end{enumerate}

\begin{enumerate}
\def\labelenumi{\alph{enumi})}
\setcounter{enumi}{1}
\tightlist
\item
  Let E be a Fréchet space. Show that if the bidual \(E''\) of E is reflexive, then E is reflexive.
\end{enumerate}

\begin{enumerate}
\def\labelenumi{\arabic{enumi})}
\setcounter{enumi}{12}
\tightlist
\item
  \begin{enumerate}
  \def\labelenumii{\alph{enumii})}
  \tightlist
  \item
    Let E, F be two Fréchet spaces, G a locally convex Hausdorff space and \(E'\) , \(F'\) , \(G'\) the duals of E, F, G respectively. Let u be a bilinear mapping from \(E' \times F'\) into \(G'\) , which is separately continuous (III, p.~28) when \(E'\) , \(F'\) , \(G'\) are assigned the weak topologies \(\sigma(E', E)\) , \(\sigma(F', F)\) and \(\sigma(G', G)\) . Show that under these conditions, u is a continuous mapping from \(E' \times F'\) into \(G'\) when \(E'\) , \(F'\) and \(G'\) are assigned the strong topologies. (For \(z \in G\) , put \(\langle z, u(x', y') \rangle = \langle v_z(x'), y' \rangle\) where \(v_z(x') \in F\) . First show that if \(E'\) is assigned the strong topology and \(F\) the initial topology, then the set of all \(v_z\) , as \(z\) ranges over a bounded set \(C\) in \(G\) , is equicontinuous; for this use IV, p.~51, exerc. 19, d). Next show that there exists a neighbourhood \(V'\) of 0 for the strong topology of \(E'\) such that the union of the sets \(v_z(V')\) as \(z\) ranges over \(C\) is bounded in \(F\) ; for this use III, p.~47, exerc. 5.)
  \end{enumerate}
\end{enumerate}

\begin{enumerate}
\def\labelenumi{\alph{enumi})}
\setcounter{enumi}{1}
\tightlist
\item
  Give an example to show that the conclusion of a) does not hold if we assume that E is a Fréchet space and F a strict inductive limit of Fréchet spaces (III, p.~47, exerc. 3).
\end{enumerate}

\begin{enumerate}
\def\labelenumi{\arabic{enumi})}
\setcounter{enumi}{13}
\tightlist
\item
  \begin{enumerate}
  \def\labelenumii{\alph{enumii})}
  \tightlist
  \item
    Let E be a Fréchet space, \(E'\) its dual. Show that \(E'\) , endowed with the topology of compact convergence or with a finer S-topology, is complete (cf.~III, p.~22, Remark 1). If E is not reflexive, show that \(E'\) is not infrabarrelled for any S-topology which is finer than the topology of compact convergence and coarser than \(\tau(\mathrm{E}', \mathrm{E})\) .
  \end{enumerate}
\end{enumerate}

\begin{enumerate}
\def\labelenumi{\alph{enumi})}
\setcounter{enumi}{1}
\tightlist
\item
  Let \((\mathrm{E}_{\alpha})_{\alpha\in\mathbf{A}}\) be a family of Fréchet spaces, E a vector space and for every \(\alpha\in A\) , let \(h_{\alpha}\) be a linear mapping from \(E_{\alpha}\) into E. Suppose that E, endowed with the finest locally convex topology for which the \(h_{\alpha}\) are continuous (II, p.~27) is Hausdorff. Prove that the dual \(E'\) of E, endowed with the topology of compact convergence or with any finer S-topology, is complete (cf.~III, p.~20, th. 1).
\end{enumerate}

\begin{enumerate}
\def\labelenumi{\arabic{enumi})}
\setcounter{enumi}{14}
\tightlist
\item
  Let \(\mathrm{E}\) be an infinite dimensional Banach space, \((a_{n})_{n\geqslant 1}\) a countable total free family of points of \(\mathrm{E}\) , and let \(\mathrm{F}_n\) be the \(n\) -dimensional subspace of \(\mathrm{E}\) generated by the \(a_{m}\) for \(m\leqslant n\) . Let \(\mathrm{S}_n\) be the sphere with equation \(\| x\| = n\) in \(\mathrm{E}\) ; in \(\mathrm{S}_n\cap \mathrm{F}_n\) let \(\mathrm{A}_n\) be a finite set such that every
\end{enumerate}

point of \(S_{n} \cap F_{n}\) is at a distance \(\leqslant 1/n\) from \(A_{n}\) . Prove that \(A = \bigcap_{n=1}^{\infty} A_{n}\) is such that its inter-

section with every closed bounded set is closed, but that 0 is a limit point of A for the weakened topology.

\begin{enumerate}
\def\labelenumi{\arabic{enumi})}
\setcounter{enumi}{15}
\tightlist
\item
  Let E be an inductive limit space of a sequence \((\mathrm{E}_{p})\) of locally convex metrizable spaces, the canonical mappings \(E_{p} \rightarrow E\) being injective and let E be the union of the images of the \(E_{p}\) . Show that the strong dual \(E_{b}^{\prime}\) of E is exhaustible (III, p.~49, exerc. 1). (Let \((\mathrm{U}_{j}^{p})\) be a decreasing fundamental system of closed convex and balanced neighbourhoods of 0 in \(E_{p}\) ; consider the finite intersections of the polar sets \((\mathrm{U}_{j}^{p})^{\circ}\) in \(E^{\prime}\) ; use the fact that for every increasing sequence
\end{enumerate}

\((m_{j})_{j\geqslant 1}\) the intersection \(\bigcap_{j = 1}^{\infty}(\mathrm{U}_{m_j}^j)^{\circ}\) is the polar of a neighbourhood of 0 in E, and that \(\mathrm{E}_b^\prime\) is complete.)

¶ 17) a) Let E be a locally convex Hausdorff space, such that the bornology consisting of the relatively compact sets of E has a countable base (A \(_{n}\) ) (III, p.~1). Show that (A \(_{n}\) ) is also a base for the canonical bornology (III, p.~3, def. 5). (Let C \(_{n}\) be the relatively compact set which is the sum of n sets each equal to A \(_{n}\) ; then (C \(_{n}\) ) is also a base for the bornology of relatively compact sets. Argue by reductio ad absurdum, considering a bounded set B which is not contained in any of the C \(_{n}\) , and conclude that there exists a sequence (x \(_{n}\) ) of points of B such that x \(_{n}\) /n∉ A \(_{n}\) ; deduce a contradiction.) Then, the space E is semi-reflexive and the closed convex balanced envelope of every compact set in E is compact.

\begin{enumerate}
\def\labelenumi{\alph{enumi})}
\setcounter{enumi}{1}
\tightlist
\item
  Suppose that E is infrabarrelled and that for a topology T compatible with the duality between E and \(E'\) , there exists a countable base for the bornology of relatively compact subsets of E for T. Show that then E is a reflexive (DF) space (IV, p.~57, exerc. 2). (Use a), observing that the closed bounded sets in E are compact for T, and consequently, complete for the initial topology of E; this implies that E is barrelled). If T is the initial topology, then E is a Montel space.
\end{enumerate}

\begin{enumerate}
\def\labelenumi{\arabic{enumi})}
\setcounter{enumi}{17}
\tightlist
\item
  \begin{enumerate}
  \def\labelenumii{\alph{enumii})}
  \tightlist
  \item
    Let E be a locally convex Hausdorff space, and let \((\mathbf{A}_{n})\) be an increasing sequence of convex, balanced, compact sets for the weakened topology \(\sigma(\mathrm{E}, \mathrm{E}')\) , such that the union of the sets \(A_{n}\) is E, and that for every integer n and every \(\lambda > 0\) , there exists m such that \(\lambda A_{n} \subset A_{m}\) . Show that \((\mathbf{A}_{n})\) is a base for the bornology consisting of convex and relatively compact sets for \(\sigma(\mathrm{E}, \mathrm{E}')\) . (Observe that on \(E'\) the topology of uniform convergence on the \(A_{n}\) is \(\tau(\mathrm{E}', \mathrm{E})\) .) b) If E is barrelled and if there exists a sequence \((\mathbf{A}_{n})\) having the properties mentioned in a), then E is a reflexive (DF) space. (Observe that \(E'\) , endowed with \(\tau(\mathrm{E}', \mathrm{E})\) is metrizable and that the topology of E is \(\beta(\mathrm{E}, \mathrm{E}')\) .)
  \end{enumerate}
\end{enumerate}

\exercisegroup

\begin{enumerate}
\def\labelenumi{\arabic{enumi})}
\item
  Let E and F be two locally convex Hausdorff spaces, \(E'\) and \(F'\) their respective duals. Show that if, for every vector subspace N of F, the topology induced on N by \(\tau(\mathrm{F}, \mathrm{F}')\) is identical with \(\tau(\mathrm{N}, \mathrm{F}'/\mathrm{N}^{\circ})\) (IV, p.~51, exerc. 20), then every strict morphism from E into F for the topologies \(\sigma(\mathrm{E}, \mathrm{E}')\) and \(\sigma(\mathrm{F}, \mathrm{F}')\) is also a strict morphism for the topologies \(\tau(\mathrm{E}, \mathrm{E}')\) and \(\tau(\mathrm{F}, \mathrm{F}')\) (cf.~IV, p.~10, prop. 11). Consider the case when F is metrizable.
\item
  Let E be a locally convex Hausdorff space such that in the dual \(E'\) there exists an infinite dimensional convex set \(B'\) which is compact for \(\sigma(E', E)\) (this condition is realized for example, when E is an infinite dimensional vector space, endowed with \(\sigma(E, E^{*})\) ). Show that there exists a linear form \(u \in E'^{*}\) which is unbounded on \(B'\) ; deduce that \(B'\) is not compact for the topology \(\sigma(E', F)\) , where F is the subspace \(E + Ru\) of \(E'^{*}\) . Conclude from this that the canonical injection from E into F is a strict morphism for the topologies \(\sigma(E, E')\) and \(\sigma(F, E')\) , but not for the topologies \(\tau(E, E')\) and \(\tau(F, E')\) .
\item
  Give an example of a strict injective morphism \(u\) from a Fréchet space \(E\) into a Fréchet space \(F\) such that \(^t u\) is not a strict morphism from \(F_b'\) into \(E_b'\) (cf.~IV, p.~58, exerc. 5, c)).
\item
  Let E and F be two locally convex Hausdorff spaces, u a continuous linear mapping from E into F, and M an everywhere dense vector subspace of E. Show that if the restriction of u to M is a strict morphism from M into F, then u is a strict morphism from E into F (use prop. 2 of IV, p.~27). If in addition \(u(\mathbf{M}) = \mathbf{F}\) , show that for every open, convex and balanced neighbourhood V of 0 in E, \(u(\mathbf{V})\) is in the interior of \(\overline{u(\mathbf{V} \cap \mathbf{M})}\) .
\item
  Let \(\mathbf{E}\) and \(\mathbf{F}\) be two normed spaces, \(u\) a continuous linear mapping from \(\mathbf{E}\) into \(\mathbf{F}\) .
\end{enumerate}

\begin{enumerate}
\def\labelenumi{\alph{enumi})}
\item
  Show that if \(u\) is a strict morphism from \(E\) into \(F\) , then \({}^t u\) is a strict morphism from the strong dual \(F_b'\) into the strong dual \(E_b'\) .
\item
  Suppose E is complete; show that if \(^{t}u\) is a strict morphism from \(F_{b}^{\prime}\) into \(E_{b}^{\prime}\) , then u is a strict morphism from E into F, and \(^{t}u\) is a strict morphism from \(F^{\prime}\) into \(E^{\prime}\) for the weak topologies \(\sigma(F^{\prime}, F)\) and \(\sigma(E^{\prime}, E)\) (consider F as a subspace of its completion).
\item
  Give an example where E is not complete, F is complete, \(^{t}u\) is a strict injective morphism from \(F'\) into \(E'\) for the strong topologies and for the weak topologies, but u is not a strict morphism from E into F (cf.~II, p.~74, exerc. 5).
\item
  Give an example where E is not complete, F is complete, u is a strict injective morphism from E into F, \(^{t}u\) is a strict morphism from F' into E' for the strong topologies, but not for the weak topologies (take E to be everywhere dense in F).
\item
  If F is complete and if \({}^t u\) is a strict morphism from \(\mathrm{F}'\) into \(\mathrm{E}'\) for the weak topologies, then \({}^t u\) is a strict morphism from \(\mathrm{F}'\) into \(\mathrm{E}'\) for the strong topologies (extend \(u\) to \(\hat{\mathrm{E}}\) ).
\item
  Give an example where E is complete, F is not complete, \(^{t}u\) is a strict morphism from \(F'\) into \(E'\) for the weak topologies but not for the strong topologies (cf.~II, p.~74, exerc. 5).
\end{enumerate}

\begin{enumerate}
\def\labelenumi{\arabic{enumi})}
\setcounter{enumi}{5}
\item
  Let E be the locally convex metrizable space \(\ell^{1}(\mathbf{N})\) (I, p.~4) endowed with the topology induced by that of the product space \(R^{N}\) ; its dual \(E'\) can be identified with \(\mathbf{R}^{(N)}\) and the topology \(\tau(\mathrm{E}',\mathrm{E})\) is the topology induced on \(E'\) by the norm topology of \(c_{0}(\mathbf{N})\) (IV, p.~47, exerc. 1). Show that if u is a surjective continuous linear mapping from E onto a Hausdorff barrelled space F, then F is necessarily finite dimensional. (Observe that \({}^{t}u\) is an isomorphism from \(F'\) , endowed with \(\sigma(\mathrm{F}',\mathrm{F})\) , onto a subspace of \(E'\) , endowed with \(\sigma(\mathrm{E}',\mathrm{E})\) ; from the fact that F is barrelled, conclude that \({}^{t}u(F')\) endowed with the topology induced by \(\tau(\mathrm{E}',\mathrm{E})\) , is a Banach space, and use exerc. 24 of II, p.~80.)
\item
  \begin{enumerate}
  \def\labelenumii{\alph{enumii})}
  \tightlist
  \item
    Let \(\mathbf{E}\) be a Banach space, and \((x_{\alpha})_{\alpha \in \mathbf{A}}\) an everywhere dense set in the unit sphere of \(\mathbf{E}\) . Let \(u\) be the linear mapping from the space \(\ell^1 (\mathbf{A})(\mathrm{I},\mathrm{p.~4})\) into \(\mathbf{E}\) defined by \(u(t) = \sum_{\alpha \in \mathbf{A}}t(\alpha)x_{\alpha}\) .
  \end{enumerate}
\end{enumerate}

for all \(t = (t(\alpha))_{\alpha \in \mathbf{A}}\) belonging to \(\ell^1 (\mathbf{A})\) . Show that \(u\) is a strict morphism from \(\ell^1 (\mathbf{A})\) onto \(\mathbf{E}\) and consequently that \(\mathbf{E}\) is isomorphic to a quotient space of \(\ell^1 (\mathbf{A})\) .

\begin{enumerate}
\def\labelenumi{\alph{enumi})}
\setcounter{enumi}{1}
\tightlist
\item
  From a) deduce an example of a closed subspace of \(\ell^{1}(\mathbf{N})\) which has no topological complement in \(\ell^{1}(\mathbf{N})\) (for E take \(c_{0}(\mathbf{N})\) and use IV, p.~54, exerc. 15, b) and p.~55, exerc. 18).
\end{enumerate}

\begin{enumerate}
\def\labelenumi{\arabic{enumi})}
\setcounter{enumi}{7}
\tightlist
\item
  For every integer \(n > 0\) , let \(a^{(n)}\) be the double sequence \(a^{(n)} = (a_{ij}^{(n)})_{i\geqslant 1,j\geqslant 1}\) defined by \(a_{ij}^{(n)} = j^n\) for every pair \((i,j)\) such that \(i < n\) and \(a_{ij}^{(n)} = i^n\) for every pair \((i,j)\) such that \(i\geqslant n\) . Let \(E\) be the vector space of double sequences \(x = (x_{ij})\) of real numbers such that, for every integer \(n > 0\) , we have \(p_n(x) = \sum_{i,j}a_{ij}^{(n)}|x_{ij}| < +\infty\) ; the semi-norms \(p_n\) define a topology of
\end{enumerate}

a Fréchet space and of a Montel space on E (IV, p.~60, exerc. 11); the dual E' of E, which is a (DF) space and a Montel space (hence ultrabornological and reflexive) is identical with the space of all sequences \(x' = (x'_{ij})\) such that, for at least one index n, we have the relation \(\sup |a_{ij}^{(n)}|^{-1} |x'_{ij}| < +\infty\) (IV, p.~47, exerc. 1, c)).

\begin{enumerate}
\def\labelenumi{\alph{enumi})}
\item
  For every \(x = (x_{ij}) \in E\) , let \(y_j = \sum_i x_{ij}\) (for all \(j \geqslant 1\) ). Show that \(\sum_j |y_j| < +\infty\) ; let \(u(x)\) denote the sequence \((y_i) \in \ell^1(\mathbf{N})\) ; show that u is a continuous linear mapping from E into \(F = \ell^1(\mathbf{N})\) , and that for every sequence \(y' = (y'_j) \in F' = \ell^\infty(\mathbf{N})\) , \(^t u(y')\) is the sequence \((z'_{ij}) \in E'\) for which \(z'_{ij} = y'_j\) for every index i. Deduce that \(^t u\) is an injective linear mapping from \(\ell^\infty(\mathbf{N})\) onto a subspace of \(E'\) which is closed for \(\sigma(E', E)\) , and consequently that u is a strict morphism from E onto \(\ell^1(\mathbf{N})\) for the initial topologies, and \(^t u\) is an isomorphism from \(F' = \ell^\infty(\mathbf{N})\) onto \(^t u(F')\) for the weak topologies \(\sigma(F', F)\) and \(\sigma(E', E)\) .
\item
  Let \(M = u^{-1}(0)\) ; then \(M^{\circ} = {}^{t}u(F')\) . Show that the inverse image under \({}^{t}u\) of the topology induced on \(M^{\circ}\) by the strong topology of \(E'\) , is the topology of uniform convergence on compact subsets of F (IV, p.~28, th. 1 and p.~54, exerc. 15). Deduce that on \(M^{\circ}\) , the topology induced by the strong topology \(\beta(E', E)\) is not the strong topology \(\beta(M^{\circ}, E/M)\) , and that for the topology induced by \(\beta(E', E)\) , \(M^{\circ}\) is not an infra-barrelled space, in spite of being a closed subspace of an ultrabornological Montel space; on the other hand, E/M, which is the quotient of a Fréchet and a Montel space by a closed subspace, is not reflexive. Show that in E/M there exist bounded sets which are not the canonical images of bounded subsets of E.
\end{enumerate}

\begin{enumerate}
\def\labelenumi{\arabic{enumi})}
\setcounter{enumi}{8}
\tightlist
\item
  Let A be a countable set. Consider three pairs of vector spaces \((P, P')\) , \((Q, Q')\) , \((E, E')\) , each of the six spaces being vector subspaces of \(R^{A}\) and containing the direct sum subspace \(\mathbf{R}^{(A)}\) ; in addition, suppose that for every point \(x \in P\) (resp. \(x \in Q\) , \(x \in E\) ) and every point \(x' \in P'\) (resp. \(x' \in Q'\) , \(x' \in E'\) ), the family \((x_{\alpha}x_{\alpha}')_{\alpha \in A}\) is summable and put \(\langle x, x' \rangle = \sum_{\alpha \in A} x_{\alpha}x_{\alpha}'\) ;
\end{enumerate}

this bilinear form puts P and P' (resp. Q and Q', E and E') in separating duality.

\begin{enumerate}
\def\labelenumi{\alph{enumi})}
\item
  Suppose that \(\mathrm{E} = \mathrm{P} \cap \mathrm{Q}\) , \(\mathrm{E}' \supset \mathrm{P}' + \mathrm{Q}'\) and \(\mathrm{E}' \neq \mathrm{P}' + \mathrm{Q}'\) . Consider the linear mapping \(u: x \mapsto (x, x)\) from \(\mathrm{E}\) into \(\mathrm{F} = \mathrm{P} \times \mathrm{Q}\) , which is put into separating duality with \(\mathrm{F}' = \mathrm{P}' \times \mathrm{Q}'\) . Show that \(u\) is continuous for the weak topologies \(\sigma(\mathrm{E}, \mathrm{E}')\) and \(\sigma(\mathrm{F}, \mathrm{F}')\) and that its image \(\mathbf{M} = u(\mathbf{E})\) is a closed subspace for \(\sigma(\mathbf{F}, \mathbf{F}')\) ; deduce that \(^t u\) is a strict morphism from \(\mathbf{F}'\) into \(\mathbf{E}'\) for the topologies \(\sigma(\mathbf{F}', \mathbf{F})\) and \(\sigma(\mathbf{E}', \mathbf{E})\) , and that \(\mathbf{N} = {}^{t}u(\mathbf{F}')\) is not closed in \(\mathbf{E}'\) for \(\sigma(\mathbf{E}', \mathbf{E})\) . If \(\mathbf{E}'\) is metrizable for the topology \(\tau(\mathbf{E}', \mathbf{E})\) , deduce that \(^t u\) is also a strict morphism from \(\mathbf{F}'\) into \(\mathbf{E}'\) for the topologies \(\tau(\mathbf{F}', \mathbf{F})\) and \(\tau(\mathbf{E}', \mathbf{E})\) .
\item
  In addition, suppose that \(E'\) is a Fréchet space for \(\tau(E', E)\) . Then \(F'/M^{\circ}\) endowed with the quotient topology of \(\tau(F', F)\) by \(M^{\circ}\) , is not semi-complete, and there exist bounded sets in \(F'/M^{\circ}\) which are not relatively compact for \(\tau(F'/M^{\circ}, M)\) .
\item
  Under the same hypotheses as in b), let \(x'\) be an element of \(E'\) not belonging to \(\mathbf{N} = {}^{u}(\mathbf{F}')\) ; for every \(y \in M\) , let \(v(y) = \langle x, x' \rangle\) , where \(x \in E\) is the unique element such that \(u(x) = y\) . Show that the linear form v on M is not continuous for the topology \(\sigma(F, F')\) , but that its restriction to every bounded subset of M is continuous for \(\sigma(M, F'/M^{\circ})\) . Deduce that \(L = v^{-1}(0)\) is a vector subspace of F, whose intersection with every bounded and closed subset of F (for \(\sigma(F, F')\) ) is closed for \(\sigma(F, F')\) , but which is not closed in F for \(\sigma(F, F')\) .
\end{enumerate}

¶ 10) a) Let G, H be two reflexive Banach spaces such that \(\mathbf{R}^{(\mathbf{N})} \subset \mathbf{G} \subset \mathbf{H} \subset \mathbf{R}^{\mathbf{N}}\) (* for example \(\mathbf{G} = \ell^{r}(\mathbf{N})\) and \(\mathbf{H} = \ell^{p}(\mathbf{N})\) , with \(1 < r < p < +\infty\) ). Taking \(\mathbf{A} = \mathbf{N} \times \mathbf{N}\) , \(\mathbf{P} = \mathbf{H}^{(\mathbf{N})}\) (topological direct sum), \(\mathbf{P}' = \mathbf{H}'^{\mathbf{N}}\) , \(\mathbf{Q} = \mathbf{G}^{\mathbf{N}}\) , \(\mathbf{Q}' = \mathbf{G}'^{(\mathbf{N})}\) , \(\mathbf{E} = \mathbf{G}^{(\mathbf{N})}\) , \(\mathbf{E}' = \mathbf{G}'^{\mathbf{N}}\) , show that the conditions of exerc. 9, a) are satisfied; that \(\mathbf{E}'\) is a reflexive Fréchet space, that E, F, F' are the strict inductive limits of reflexive Fréchet spaces (hence complete and reflexive).

\begin{enumerate}
\def\labelenumi{\alph{enumi})}
\setcounter{enumi}{1}
\tightlist
\item
  From a) and exerc. 9, construct examples with the following properties :
\end{enumerate}

\(\alpha)\) A quotient space of a strict inductive limit of reflexive Fréchet spaces which is neither quasi-complete nor semi-reflexive.

β) A closed subspace of a strict inductive limit of reflexive Fréchet spaces which is not reflexive, and whose dual is not complete.

\(\gamma)\) A vector subspace of a strict inductive limit of reflexive Fréchet spaces \(\mathbf{E}_n\) , which is not closed (hence not barrelled) but whose intersection with each of the subspaces \(\mathbf{E}_n\) is closed.

δ) A non-closed subspace of the dual of a strict inductive limit of reflexive Fréchet spaces, whose intersection with every weakly compact subset is weakly compact (cf.~IV, p.~25, cor. 2).

\[
\mathrm{E} _ {n}
\]

\[
\mathrm{F} _ {n} (c f.
\]

\[
\mathrm{F} _ {n} = \overline {{\mathrm{F} \cap \mathrm{E} _ {n}}}
\]

\[
u _ {n}
\]

\[
v _ {n}
\]

\[
\mathrm{F} \cap \mathrm{E} _ {n}
\]

\[
\overline {{\mathrm{F} \cap \mathrm{E} _ {n}}}
\]

\[
v _ {n}
\]

\begin{enumerate}
\def\labelenumi{\alph{enumi})}
\setcounter{enumi}{1}
\tightlist
\item
  Suppose F is not closed in E but that \(F \cap E_{n}\) is closed in \(E_{n}\) for all n (exerc. 10, b)). Let \(A_{n}\) be a countable dense set in \(E_{n}\) , and let G be the vector subspace of E generated by the union of F and the \(A_{n}\) . Show that G is a bornological space in which F is a subspace with a complement which has a countable basis, but that F is not infra-barrelled (cf.~III, p.~41, § 2, exerc. 4 and p.~45, exerc. 17).
\end{enumerate}

\begin{enumerate}
\def\labelenumi{\arabic{enumi})}
\setcounter{enumi}{11}
\tightlist
\item
  Let E, F be two Fréchet spaces, u a strict morphism from E into F. Show that for every finite rank continuous linear mapping v from E into F, \(u + v\) is a strict morphism from E into F.
\end{enumerate}

¶ 13) a) Let E be a Hausdorff and complete locally convex space. Suppose that there exists ( \(F_{n}\) ) a decreasing sequence of closed vector subspaces in E such that for every neighbourhood V of 0 in E, there exists n such that \(F_{n} \subset V\) . Show that the space E is of minimal type (II, p.~85, exerc. 13).

\begin{enumerate}
\def\labelenumi{\alph{enumi})}
\setcounter{enumi}{1}
\item
  Let E be a Fréchet space satisfying the first axiom of countability, and which is not of minimal type. Show that there exist two closed vector subspaces M, N in E such that \(M \cap N = \{0\}\) and such that \(M + N\) is not closed. (Let \((x_{n}')\) denote a sequence of linearly independent continuous linear forms on E, forming a total set for \(\sigma(E', E)\) (III, p.~19, cor. 2); let \(L_{n}\) be the subspace of E orthogonal to the \(x_{i}'\) for indices \(i \leqslant 2n\) ; let \(x_{n}, y_{n}\) be two linearly independent vectors in the complement of \(L_{n+1}\) with respect to \(L_{n}\) . Let d be a translation invariant distance defining the topology of E. Using a), show that there exists a number \(\alpha > 0\) such that we can take \(d(0, x_{n}) \geqslant \alpha\) , \(d(0, y_{n}) \geqslant \alpha\) and \(d(x_{n}, y_{n}) \leqslant 1/n\) . Show that if M (resp. N) is the closed vector subspace generated by the \(x_{n}\) (resp. \(y_{n}\) ), then M and N have the properties required; use I, p.~19, cor. 4.)
\item
  Let \(\mathbf{E}\) be a closed subspace of a product \(\prod_{\alpha \in \mathbf{A}} \mathrm{F}_{\alpha}\) of normed spaces; show that if every closed
\end{enumerate}

subspace of E generated by a countable family of points is of minimal type, then E itself is of minimal type. (Argue by reductio ad absurdum, assuming that the projection of E on each \(F_{\alpha}\) is equal to \(F_{\alpha}\) , and that for some \(\alpha\) , \(F_{\alpha}\) is infinite dimensional.)

\begin{enumerate}
\def\labelenumi{\alph{enumi})}
\setcounter{enumi}{3}
\tightlist
\item
  Let E be a Hausdorff and complete locally convex space, in which every subspace generated by a countable family of points is metrizable. In order that E be of minimal type, it is necessary and sufficient that for every pair of closed vector subspaces M, N of E such that \(M \cap N = \{0\}\) , \(M + N\) is closed in E (use b) and c)).
\end{enumerate}

\begin{enumerate}
\def\labelenumi{\arabic{enumi})}
\setcounter{enumi}{13}
\tightlist
\item
  \begin{enumerate}
  \def\labelenumii{\alph{enumii})}
  \tightlist
  \item
    Let E be a Fréchet space on which there exists no continuous norm. Show that there exists a closed vector subspace of minimal type in E (II, p.~85, exerc. 13), which is infinite dimensional, and consequently, has a topological complement in E. (There exists a strictly decreasing fundamental sequence \((\mathrm{V}_{n})\) of convex, balanced neighbourhoods of 0 in E, and a sequence \((x_{n})\) of points of E such that \(x_{n} \notin V_{n+1}\) and such that the line passing through \(x_{n}\) is contained in \(V_{n}\) , the \(x_{n}\) being linearly independent; show that the required space is the closed vector subspace generated by the \(x_{n}\) .)
  \end{enumerate}
\end{enumerate}

\begin{enumerate}
\def\labelenumi{\alph{enumi})}
\setcounter{enumi}{1}
\item
  Let E be a Fréchet space whose topology cannot be defined by a single norm, but by an increasing sequence \((p_{n})\) of norms. Let \(V_{n}\) be the neighbourhood defined by \(p_{n}(x) \leqslant 1\) , and let \(A_{n}^{\prime} = V_{n}^{\circ}\) be its polar in \(E^{\prime}\) ; we can assume that \(A_{n+1}^{\prime}\) is not contained in the vector subspace of \(E^{\prime}\) generated by \(A_{n}^{\prime}\) (IV, p.~49, exerc. 12, c)). Let \((x_{n}^{\prime})\) be a sequence of points of \(E^{\prime}\) such that \(x_{n}^{\prime} \in A_{n}^{\prime}\) and that \(x_{n}^{\prime}\) does not belong to the vector subspace generated by \(A_{n-1}^{\prime}\) ; show that the vector subspace \(M^{\prime}\) generated by the sequence \((x_{n}^{\prime})\) is weakly closed in \(E^{\prime}\) (cf.~IV, p.~25, cor. 2) and does not have a topological complement in \(E^{\prime}\) for the topology \(\sigma(E^{\prime}, E)\) . (Observe that if there existed a weakly continuous projector \(u^{\prime}\) from \(E^{\prime}\) onto \(M^{\prime}\) , then \(u^{\prime}(A_{1}^{\prime})\) would be contained in one of the sets \(M^{\prime} \cap A_{n}^{\prime}\) by Baire's theorem, and derive a contradiction, since \(A_{1}^{\prime}\) is weakly total in \(E^{\prime}\) .) Deduce that the subspace \(M^{\circ}\) of E does not have a topological complement in E.
\item
  Let E be a Fréchet space whose topology cannot be defined by a single norm. Show that if E is not isomorphic to a product of Banach spaces, then there exists a closed vector subspace in E which has no topological complement. (Argue by reductio ad absurdum; let \((p_{n})_{n\geqslant1}\) be an increasing sequence of semi-norms on E, defining the topology of E; let \(F_{n}=p_{n}^{-1}(0)\) and let \(E_{n+1}\) be a topological complement of \(F_{n+1}\) with respect to \(F_{n}\) (we put \(E=F_{0}\) ); using b), show that \(E_{n+1}\) is a Banach space and that E is isomorphic to the product of the \(E_{n}(n\geqslant1)\) ; for this use I, p.~17, th. 1.)
\end{enumerate}

\begin{enumerate}
\def\labelenumi{\arabic{enumi})}
\setcounter{enumi}{14}
\tightlist
\item
  \begin{enumerate}
  \def\labelenumii{\alph{enumii})}
  \tightlist
  \item
    Let E be an infinite dimensional normed space (real or complex). Show that there exists a sequence \((x_{n})\) in E such that, for every bounded sequence \((\lambda_{n})\) of scalars, there exists a continuous linear form \(x^{\prime}\) on E such that \(\langle x_n,x'\rangle = \lambda_n\) for all \(n\) . (Construct a sequence \((x_n')\) of points in the dual E' of E and a sequence \((x_{n})\) of points in E such that \(\langle x_i,x_j'\rangle = \delta_{ij}\) and \(\| x_n'\| \leqslant 2^{-n}\) .)
  \end{enumerate}
\end{enumerate}

\begin{enumerate}
\def\labelenumi{\alph{enumi})}
\setcounter{enumi}{1}
\tightlist
\item
  For a locally convex metrizable space \(\mathbf{E}\) to have the property stated in \(a\) ), it is necessary and sufficient that the completion of \(\mathbf{E}\) is not a space of minimal type (II, p.~85, exerc. 13).
\end{enumerate}

\begin{enumerate}
\def\labelenumi{\arabic{enumi})}
\setcounter{enumi}{15}
\tightlist
\item
  \begin{enumerate}
  \def\labelenumii{\alph{enumii})}
  \tightlist
  \item
    Let E and F be two infinite dimensional complex normed spaces, and u a bijective semilinear mapping from E onto F, with respect to an automorphism \(\sigma\) of C, which transforms every closed hyperplane of E into a closed hyperplane of F. Show that the automorphism \(\sigma\) of C is necessarily continuous (and consequently, is either the identity or the automorphism \(\xi \mapsto \overline{\xi}\) ). (Argue by reductio ad absurdum : let \((x_{n})\) be a sequence of points in E satisfying the condition of exerc. 15, a) and let \((\lambda_{n})\) be a bounded sequence of complex numbers such that \(|\lambda_{n}^{\sigma}| \geqslant n\) . \(\|u(x_{n})\|\) for all n; if \(x' \in E'\) is such that \(\langle x_{n}, x' \rangle = \lambda_{n}\) for all n, consider the image under u of the closed hyperplane of all \(x \in E\) such that \(\langle x, x' \rangle = 1\) and derive a contradiction.)
  \end{enumerate}
\end{enumerate}

\begin{enumerate}
\def\labelenumi{\alph{enumi})}
\setcounter{enumi}{1}
\tightlist
\item
  Deduce from a) that u is a continuous semi-linear mapping from E into F (IV, p.~7, corollary).
\item
  Let \(\sigma\) be a discontinuous automorphism of the field C. Show that the bijection \((\xi_{n})\mapsto(\xi_{n}^{\sigma})\) from \(C^{N}\) onto itself transforms every closed hyperplane into a closed hyperplane (cf.~IV, p.~14, prop. 15).
\end{enumerate}

17) Let E, F be two Banach spaces, u a strict morphism from E onto F. Then there exists
a number \(m > 0\) such that for every \(\varepsilon \in ]0, m[\) and for every \(y \in F\), there exists \(z \in E\) such
that \(u(z) = y\) and \(\| z \| \leqslant (m - \varepsilon)^{-1} \| y \|\).
a) Let \(\mathbf{B}\) be a closed ball \(\| x - a \| \leqslant r\) in E, and let w be a mapping from \(\mathbf{B}\) into F satisfying
the following conditions: \(1^{\circ}\ \sup_{x \in \mathbf{B}} \| w(x) \| = M < +\infty; 2^{\circ}\) there exists a number \(k > 0\)
such that for all \(x, x'\) in \(\mathbf{B}\), we have \(\| w(x) - w(x') \| \leqslant k \| x - x' \|\) («Lipschitz condition »). Show that if \(k < m\) and \(M < r(m - k)\), then the image of \(\mathbf{B}\) under the mapping \(x \mapsto u(x) + w(x)\)
contains a ball with center \(b = u(a)\). (Show that for every \(y \in F\) close enough to \(b\), we can define
a sequence \((x_n)_{n \geqslant 0}\) of points of \(\mathbf{B}\) such that \(u(x_0) = y\) and \(u(x_n) = y - w(x_{n-1})\) for \(n \geqslant 1\),
and such that \((x_n)\) converges to a point of \(\mathbf{B}\).)


\begin{enumerate}
\def\labelenumi{\alph{enumi})}
\setcounter{enumi}{1}
\tightlist
\item
  Suppose that \(w\) is a mapping from all \(E\) into \(F\) , such that \(\| w(x) - w(x') \| \leqslant k \| x - x' \|\) for every \(x, x'\) in \(E\) . Show that if \(k < m\) , \(u + w\) is a surjective and open mapping from \(E\) onto \(F\) . c) In particular, if \(w\) is a continuous linear mapping from \(E\) into \(F\) such that \(\| w \| < m\) , then \(v = u + w\) is a strict morphism from \(E\) onto \(F\) . If \(\varepsilon \in ]0, m[\) , then for every \(x \in v^{-1}(0)\) with \(\| x \| = 1\) , there exists \(x_0 \in u^{-1}(0)\) such that \(\| x - x_0 \| \leqslant \| w \| / (m - \varepsilon)\) ; if in addition \(\| w \| < m - \varepsilon\) , then for every \(x_0 \in u^{-1}(0)\) with \(\| x_0 \| = 1\) , there exists \(x \in v^{-1}(0)\) such that
\end{enumerate}

\[
\left\| x - x _ {0} \right\| \leqslant \| w \| / (m - \varepsilon - \| w \|).
\]

Deduce that if there exists a closed vector subspace G of E which is the complement of \(u^{-1}(0)\) , then G is also the complement of \(v^{-1}(0)\) whenever \(\|w\|\) is small enough (argue by reductio ad absurdum: show that the projection of \(v^{-1}(0)\) onto \(u^{-1}(0)\) cannot be contained in a closed hyperplane of \(u^{-1}(0)\) ; on the other hand, note that there exists a > 0 such that every point \(x \in G\) such that \(\|x\| = 1\) is at a distance \(\geqslant a\) from \(u^{-1}(0)\) ).

\begin{enumerate}
\def\labelenumi{\arabic{enumi})}
\setcounter{enumi}{17}
\tightlist
\item
  \begin{enumerate}
  \def\labelenumii{\alph{enumii})}
  \tightlist
  \item
    Let E, F be two normed spaces, u a strict injective morphism from E into F; then there is a number m > 0 such that \(\|u(x)\| \geqslant m \|x\|\) for all \(x \in E\) . Show that if w is a continuous linear mapping from E into F such that \(\|w\| < m\) , then \(v = u + w\) is a strict injective morphism from E into F. Moreover, for all \(y_0 \in u(E)\) such that \(\|y_0\| = 1\) , there exists \(y \in v(E)\) such that \(\|y - y_0\| \leqslant \|w\| / m\) ; for every \(y \in v(E)\) with \(\|y\| = 1\) , there exists \(y_0 \in u(E)\) such that \(\|y - y_0\| \leqslant \|w\| / (m - \|w\|)\) .
  \end{enumerate}
\end{enumerate}

\begin{enumerate}
\def\labelenumi{\alph{enumi})}
\setcounter{enumi}{1}
\tightlist
\item
  Deduce from a) that, if in addition, E and F are Banach spaces, and if there exists a closed vector subspace G of F which is the complement of the closed subspace \(u(\mathrm{E})\) , then G is also the complement of \(v(\mathrm{E})\) whenever \(\|w\|\) is small enough (argue as in exerc. 17, c)).
\end{enumerate}

\begin{enumerate}
\def\labelenumi{\arabic{enumi})}
\setcounter{enumi}{18}
\item
  Let E be the subspace of the Banach space \(\mathcal{C}(\left[-1,1\right];\mathbf{R})\) of all continuous mappings from \([-1,1]\) into R, consisting of the polynomials. Similarly, let F be the subspace of \(\mathcal{C}(\left[0,1\right];\mathbf{R})\) consisting of all polynomials. Let u be the mapping which associates to each polynomial \(f\in E\) , the polynomial \(t\mapsto\frac{1}{2}(f(\sqrt{t})+f(-\sqrt{t}))\) in F. Also, let w be the mapping from E into F which associates with every polynomial \(f\in E\) its restriction to \([0,1]\) . Show that u is a strict morphism from E onto F, but that for every \(\varepsilon>0\) , \(u+\varepsilon w\) is not a strict morphism from E into F.
\item
  \begin{enumerate}
  \def\labelenumii{\alph{enumii})}
  \tightlist
  \item
    Let E, F be two Banach spaces, u a strict morphism from E into F, such that \(u^{-1}(0)\) is finite dimensional. Show that for every continuous linear mapping w from E into F with small enough norm, \(v = u + w\) is a strict morphism from E into F and \(\dim v^{-1}(0) \leqslant \dim u^{-1}(0)\) . (Write E as the topological direct sum of \(u^{-1}(0)\) and a closed subspace and use exerc. 18 of IV, p.~66.)
  \end{enumerate}
\end{enumerate}

\begin{enumerate}
\def\labelenumi{\alph{enumi})}
\setcounter{enumi}{1}
\tightlist
\item
  Let E, F be two Banach spaces, u a continuous linear mapping from E into F such that \(u(\mathrm{E})\) has finite codimension in F. Then \(u(\mathrm{E})\) is closed and u is a strict morphism (I, p.~28, exerc. 4). Show that for every continuous linear mapping w from E into F, of small enough norm, \(v = u + w\) is a strict morphism from E into F and
\end{enumerate}

\[
\operatorname{codim} (v (\mathrm{E})) \leqslant \operatorname{codim} (u (\mathrm{E}))
\]

(consider \({}^t v = {}^t u + {}^t w\) , and use a) and IV, p.~30, cor. 3).

\begin{enumerate}
\def\labelenumi{\arabic{enumi})}
\setcounter{enumi}{20}
\tightlist
\item
  Let E and F be two Banach spaces. A continuous linear mapping from E into F is said to be a Fredholm operator (or a quasi-isomorphism) if \(u^{-1}(0)\) is finite dimensional and \(u(\mathrm{E})\) has finite codimension (this implies that \(u(\mathrm{E})\) is closed in F and u is a strict morphism); the number \(\operatorname{Ind}(u) = \operatorname{codim}(u(\mathrm{E})) - \dim(u^{-1}(0))\) is called the index of u.
\end{enumerate}

\begin{enumerate}
\def\labelenumi{\alph{enumi})}
\item
  Show that \({}^t u: \mathrm{F}' \to \mathrm{E}'\) is also a Fredholm operator and that \(\operatorname{Ind}({}^t u) = -\operatorname{Ind}(u)\) .
\item
  If \(u: \mathrm{E} \to \mathrm{F}\) and \(v: \mathrm{F} \to \mathrm{G}\) are two Fredholm operators, then so is \(v \circ u: \mathrm{E} \to \mathrm{G}\) and \(\operatorname{Ind}(v \circ u) = \operatorname{Ind}(u) + \operatorname{Ind}(v)\) .
\item
  If \(w: E \to F\) is a continuous linear mapping with finite rank or with a small enough norm, then \(u + w\) is a Fredholm operator and \(\operatorname{Ind}(u + w) = \operatorname{Ind}(u)\) (use exerc. 17, c) of IV, p.~65 and 18, b)).
\end{enumerate}

\begin{enumerate}
\def\labelenumi{\arabic{enumi})}
\setcounter{enumi}{21}
\tightlist
\item
  Let \(\mathbf{X}\) be a real Banach space, \(\mathbf{E}\) a finite dimensional subspace of \(\mathbf{X}\) .
\end{enumerate}

\begin{enumerate}
\def\labelenumi{\alph{enumi})}
\item
  Let S be the unit sphere in X. Show that for every \(\varepsilon > 0\) , there exists a finite number of linearly independent points \(z_{i} \in S\) ( \(1 \leqslant i \leqslant r\) ) such that for every \(x \in S \cap E\) , there exists an index \(i\) such that \(\| x - z_{i} \| \leqslant \varepsilon\) .
\item
  Let \(z_{i}^{\prime}(1 \leqslant i \leqslant r)\) be points in the dual \(E^{\prime}\) of \(E\) with \(\| z_{i}^{\prime}\| \geqslant 1\) for every \(i\) , and \(\langle z_i, z_j' \rangle = \delta_{ij}\) (Kronecker's index), and let \(F\) be the closed subspace of codimension \(r\) in \(E\) which is orthogonal to the subspace of \(\mathbf{E}'\) generated by the \(z_{i}^{\prime}\) . Show that for every \(x\in \mathrm{S}\cap \mathrm{E}\) and for every \(y\in \mathrm{F}\) , \(\| x + y\| \geqslant 1 - \varepsilon\) , and in particular, \(\mathrm{E}\cap \mathrm{F} = \{0\}\) , so that the sum \(\mathrm{E} + \mathrm{F}\) is the topological direct sum.
\item
  Deduce from \(b)\) that the continuous projector \(P\) from \(\mathrm{E} + \mathrm{F}\) onto \(\mathrm{E}\) corresponding to the topological direct sum decomposition \(\mathrm{E} \oplus \mathrm{F}\) has a norm \(\| P \| < 1 / (1 - \varepsilon)\) .
\end{enumerate}

\begin{enumerate}
\def\labelenumi{\arabic{enumi})}
\setcounter{enumi}{22}
\tightlist
\item
  Let \(X\) be a real infinite dimensional Banach space.
\end{enumerate}

\begin{enumerate}
\def\labelenumi{\alph{enumi})}
\item
  Suppose that for every \(\lambda > 0\) , there exists a finite dimensional subspace \(\mathrm{E}_{\lambda}\) of X such that there exists no continuous projector \(P\) on X with image \(\mathrm{E}_{\lambda}\) and such that \(\| P \| \leqslant \lambda\) . Show that every closed subspace Y of X with finite codimension has the same property as X.
\item
  By induction on n, show that there exists a decreasing sequence \((\mathrm{X}_{n})\) of closed subspaces of X, of finite codimension, and for every n, a finite dimensional subspace \(E_{n}\) of X such that: \(1^{\circ}\) the sum \(E_{1} + E_{2} + \cdots + E_{n}\) is direct, and there exists a continuous projector \(P_{n}\) on \(X_{n}\) , with norm \(\leqslant 2\) whose image is \(E_{1} + E_{2} + \cdots + E_{n}\) ; \(2^{\circ}\) the space \(E_{n}\) is contained in \((\mathrm{I} - P_{n-1})(\mathrm{X}_{n-1})\) ; \(3^{\circ}\) there exists no continuous projector on X with image \(E_{n}\) and with norm \(\leqslant n + 2\) .
\item
  Let \(\bar{Z}\) be the closed subspace of \(X\) generated by the union of the \(E_{n}\) . Show that there does not exist a topological complement of \(Z\) in \(X\) (observe that if \(Q\) were a continuous projector on \(X\) , with image \(\bar{Z}\) , then \((I - P_{n-1})P_nQ\) would be a continuous projector on \(X\) with image \(E_{n}\) ).
\end{enumerate}
